% ═══════════════════════════════════════════════════════════════════════
\section{Autosimilarité}
% ═══════════════════════════════════════════════════════════════════════

\begin{frame}{Des copies réduites d'elle-même}
  \small
  \begin{itemize}
    \item Figure 3.3 : le premier quart de $K$, agrandi trois fois,
      redonne $K$. Idem pour les trois autres quarts (en les faisant aussi
      tourner) : $K$ est formée de quatre copies réduites d'elle-même.
    \item Le triangle de Sierpiński : trois copies de lui-même, deux fois
      plus petites, dans ses trois coins.
    \item Langage du chapitre 1 : une similitude de rapport $r$ multiplie
      les distances par $r$ ; si $0<r<1$, c'est une contraction de
      rapport $r$ (définition 2.14).
  \end{itemize}
\end{frame}

\begin{frame}{Définition 3.1 --- Ensemble autosimilaire}
  \begin{Enonce}{Définition 3.1 --- Ensemble autosimilaire}
    Soit $E$ un sous-ensemble compact non vide du plan. $E$ est
    \defi{autosimilaire} s'il existe un nombre fini $N\geq2$ de
    similitudes $w_1,\dots,w_N$, de rapports $r_1,\dots,r_N$ compris
    strictement entre $0$ et $1$, telles que
    \[
      E = w_1(E)\cup w_2(E)\cup\cdots\cup w_N(E).
    \]
    Chaque $w_i(E)$ est une \defi{copie réduite} de $E$.
  \end{Enonce}
  \small
  \begin{itemize}
    \item $E$ doit être exactement la réunion de ses copies.
    \item Les copies peuvent se toucher, et même se chevaucher (on y
      revient en section 8).
  \end{itemize}
\end{frame}

\begin{frame}{Exemple 3.1 --- Le segment et le carré}
  \small
  \begin{itemize}
    \item $\intff{0}{1}$ est la réunion de $\intff{0}{\tfrac12}$ et de
      $\intff{\tfrac12}{1}$ : deux similitudes de rapport $\tfrac12$ (ou
      trois de rapport $\tfrac13$, ou $m$ de rapport $\tfrac1m$).
    \item Le carré unité $C$ est la réunion de quatre carrés de côté
      $\tfrac12$ : $C = w_1(C)\cup w_2(C)\cup w_3(C)\cup w_4(C)$, avec
      \[
        w_i\col{x}{y} = \mat{\tfrac12}{0}{0}{\tfrac12}\col{x}{y} + \vec t_i,
      \]
      $\vec t_i$ parcourant $(0,0)$, $\big(\tfrac12,0\big)$,
      $\big(0,\tfrac12\big)$, $\big(\tfrac12,\tfrac12\big)$.
  \end{itemize}
\end{frame}

\begin{frame}{Exemple 3.2 --- Tout triangle est autosimilaire}
  \begin{columns}[c]
    \begin{column}{0.52\textwidth}
      \small En joignant les milieux des côtés de $T$, on obtient quatre
      triangles semblables à $T$, de rapport $\tfrac12$ :
      \begin{itemize}
        \item trois copies de coin : $h_{A,1/2}(T)$, $h_{B,1/2}(T)$,
          $h_{C,1/2}(T)$ ;
        \item la copie centrale, qui a subi un demi-tour :
          $h_{G,-1/2}(T)$, où $G$ est le centre de gravité.
      \end{itemize}
      \[
        \begin{aligned}
          T = {}& h_{A,1/2}(T)\cup h_{B,1/2}(T)\\
          &\cup h_{C,1/2}(T)\cup h_{G,-1/2}(T).
        \end{aligned}
      \]
    \end{column}
    \begin{column}{0.46\textwidth}
      \FigBox[0.55\textheight]{%
        \begin{tikzpicture}[scale=1.15]
          \coordinate (A) at (0,0); \coordinate (B) at (4.2,0.3); \coordinate (C) at (1.3,2.6);
          \coordinate (MAB) at ($(A)!0.5!(B)$); \coordinate (MBC) at ($(B)!0.5!(C)$);
          \coordinate (MCA) at ($(C)!0.5!(A)$);
          \fill[firebrick3!30] (A) -- (MAB) -- (MCA) -- cycle;
          \fill[steelblue4!35] (MAB) -- (B) -- (MBC) -- cycle;
          \fill[coveraccent!45] (MCA) -- (MBC) -- (C) -- cycle;
          \fill[vertpiece!30] (MAB) -- (MBC) -- (MCA) -- cycle;
          \draw[sectioncolor, thick, line join=round] (A) -- (B) -- (C) -- cycle;
          \draw[sectioncolor, thin] (MAB) -- (MBC) -- (MCA) -- cycle;
          \coordinate (G) at ($1/3*(A)+1/3*(B)+1/3*(C)$);
          \node[pt, label={[etiq]below left:$A$}] at (A) {};
          \node[pt, label={[etiq]below right:$B$}] at (B) {};
          \node[pt, label={[etiq]above:$C$}] at (C) {};
          \node[ptim, label={[etiq]above:$G$}] at (G) {};
          \node[etiq] at ($1/3*(A)+1/3*(MAB)+1/3*(MCA)$) {$T_A$};
          \node[etiq] at ($1/3*(MAB)+1/3*(B)+1/3*(MBC)$) {$T_B$};
          \node[etiq] at ($1/3*(MCA)+1/3*(MBC)+1/3*(C)$) {$T_C$};
          \node[etiq] at ($(G)+(0,-0.3)$) {$T_G$};
        \end{tikzpicture}}
      \begin{center}\scriptsize Figure 3.6\end{center}
    \end{column}
  \end{columns}
  \small L'autosimilarité ne suffit donc pas à caractériser une fractale.
\end{frame}

\begin{frame}{Exemples 3.3 et 3.4 --- Sierpiński et von Koch}
  \small
  Sierpiński, sommets $S_1=(0,0)$, $S_2=(1,0)$,
  $S_3=\big(\tfrac12,\tfrac{\sqrt3}2\big)$, $w_i=h_{S_i,1/2}$ :
  \[
    w_1\col{x}{y} = \tfrac12\col{x}{y},\quad
    w_2\col{x}{y} = \tfrac12\col{x}{y}+\col{\tfrac12}{0},\quad
    w_3\col{x}{y} = \tfrac12\col{x}{y}+\col{\tfrac14}{\tfrac{\sqrt3}{4}}.
  \]
  Von Koch, sur le segment de $(0,0)$ à $(1,0)$, avec les rotations
  $r_{\pm\pi/3}$ :
  \[
    \begin{aligned}
      w_1\col{x}{y} &= \tfrac13\col{x}{y}, &
      w_2\col{x}{y} &= \tfrac13\,r_{\pi/3}\col{x}{y}+\col{\tfrac13}{0},\\
      w_3\col{x}{y} &= \tfrac13\,r_{-\pi/3}\col{x}{y}+\col{\tfrac12}{\tfrac{\sqrt3}{6}}, &
      w_4\col{x}{y} &= \tfrac13\col{x}{y}+\col{\tfrac23}{0}.
    \end{aligned}
  \]
  Vérification : $w_2$ envoie $(1,0)$ sur
  $\big(\tfrac12,\tfrac{\sqrt3}{6}\big)$, le sommet de la pointe. Sierpiński
  est aussi formé de neuf copies de rapport $\tfrac14$, les
  $w_i\circ w_j$.
\end{frame}

\begin{frame}{Figure 3.7 --- Les copies réduites}
  \FigBox[0.58\textheight]{%
    \begin{tikzpicture}[scale=4.3]
      % Triangle de Sierpiński : les trois copies w_i(S)
% Fichier généré par generer.py : ne pas modifier à la main.
\fill[piece1]
  (0,0) -- (0.0156,0) -- (0.0078,0.0135) -- cycle (0.0156,0) -- (0.0312,0) -- (0.0234,0.0135) -- cycle (0.0078,0.0135) -- (0.0234,0.0135) -- (0.0156,0.0271) -- cycle (0.0312,0) -- (0.0469,0) -- (0.0391,0.0135) -- cycle
  (0.0469,0) -- (0.0625,0) -- (0.0547,0.0135) -- cycle (0.0391,0.0135) -- (0.0547,0.0135) -- (0.0469,0.0271) -- cycle (0.0156,0.0271) -- (0.0312,0.0271) -- (0.0234,0.0406) -- cycle (0.0312,0.0271) -- (0.0469,0.0271) -- (0.0391,0.0406) -- cycle
  (0.0234,0.0406) -- (0.0391,0.0406) -- (0.0312,0.0541) -- cycle (0.0625,0) -- (0.0781,0) -- (0.0703,0.0135) -- cycle (0.0781,0) -- (0.0938,0) -- (0.0859,0.0135) -- cycle (0.0703,0.0135) -- (0.0859,0.0135) -- (0.0781,0.0271) -- cycle
  (0.0938,0) -- (0.1094,0) -- (0.1016,0.0135) -- cycle (0.1094,0) -- (0.125,0) -- (0.1172,0.0135) -- cycle (0.1016,0.0135) -- (0.1172,0.0135) -- (0.1094,0.0271) -- cycle (0.0781,0.0271) -- (0.0938,0.0271) -- (0.0859,0.0406) -- cycle
  (0.0938,0.0271) -- (0.1094,0.0271) -- (0.1016,0.0406) -- cycle (0.0859,0.0406) -- (0.1016,0.0406) -- (0.0938,0.0541) -- cycle (0.0312,0.0541) -- (0.0469,0.0541) -- (0.0391,0.0677) -- cycle (0.0469,0.0541) -- (0.0625,0.0541) -- (0.0547,0.0677) -- cycle
  (0.0391,0.0677) -- (0.0547,0.0677) -- (0.0469,0.0812) -- cycle (0.0625,0.0541) -- (0.0781,0.0541) -- (0.0703,0.0677) -- cycle (0.0781,0.0541) -- (0.0938,0.0541) -- (0.0859,0.0677) -- cycle (0.0703,0.0677) -- (0.0859,0.0677) -- (0.0781,0.0812) -- cycle
  (0.0469,0.0812) -- (0.0625,0.0812) -- (0.0547,0.0947) -- cycle (0.0625,0.0812) -- (0.0781,0.0812) -- (0.0703,0.0947) -- cycle (0.0547,0.0947) -- (0.0703,0.0947) -- (0.0625,0.1083) -- cycle (0.125,0) -- (0.1406,0) -- (0.1328,0.0135) -- cycle
  (0.1406,0) -- (0.1562,0) -- (0.1484,0.0135) -- cycle (0.1328,0.0135) -- (0.1484,0.0135) -- (0.1406,0.0271) -- cycle (0.1562,0) -- (0.1719,0) -- (0.1641,0.0135) -- cycle (0.1719,0) -- (0.1875,0) -- (0.1797,0.0135) -- cycle
  (0.1641,0.0135) -- (0.1797,0.0135) -- (0.1719,0.0271) -- cycle (0.1406,0.0271) -- (0.1562,0.0271) -- (0.1484,0.0406) -- cycle (0.1562,0.0271) -- (0.1719,0.0271) -- (0.1641,0.0406) -- cycle (0.1484,0.0406) -- (0.1641,0.0406) -- (0.1562,0.0541) -- cycle
  (0.1875,0) -- (0.2031,0) -- (0.1953,0.0135) -- cycle (0.2031,0) -- (0.2188,0) -- (0.2109,0.0135) -- cycle (0.1953,0.0135) -- (0.2109,0.0135) -- (0.2031,0.0271) -- cycle (0.2188,0) -- (0.2344,0) -- (0.2266,0.0135) -- cycle
  (0.2344,0) -- (0.25,0) -- (0.2422,0.0135) -- cycle (0.2266,0.0135) -- (0.2422,0.0135) -- (0.2344,0.0271) -- cycle (0.2031,0.0271) -- (0.2188,0.0271) -- (0.2109,0.0406) -- cycle (0.2188,0.0271) -- (0.2344,0.0271) -- (0.2266,0.0406) -- cycle
  (0.2109,0.0406) -- (0.2266,0.0406) -- (0.2188,0.0541) -- cycle (0.1562,0.0541) -- (0.1719,0.0541) -- (0.1641,0.0677) -- cycle (0.1719,0.0541) -- (0.1875,0.0541) -- (0.1797,0.0677) -- cycle (0.1641,0.0677) -- (0.1797,0.0677) -- (0.1719,0.0812) -- cycle
  (0.1875,0.0541) -- (0.2031,0.0541) -- (0.1953,0.0677) -- cycle (0.2031,0.0541) -- (0.2188,0.0541) -- (0.2109,0.0677) -- cycle (0.1953,0.0677) -- (0.2109,0.0677) -- (0.2031,0.0812) -- cycle (0.1719,0.0812) -- (0.1875,0.0812) -- (0.1797,0.0947) -- cycle
  (0.1875,0.0812) -- (0.2031,0.0812) -- (0.1953,0.0947) -- cycle (0.1797,0.0947) -- (0.1953,0.0947) -- (0.1875,0.1083) -- cycle (0.0625,0.1083) -- (0.0781,0.1083) -- (0.0703,0.1218) -- cycle (0.0781,0.1083) -- (0.0938,0.1083) -- (0.0859,0.1218) -- cycle
  (0.0703,0.1218) -- (0.0859,0.1218) -- (0.0781,0.1353) -- cycle (0.0938,0.1083) -- (0.1094,0.1083) -- (0.1016,0.1218) -- cycle (0.1094,0.1083) -- (0.125,0.1083) -- (0.1172,0.1218) -- cycle (0.1016,0.1218) -- (0.1172,0.1218) -- (0.1094,0.1353) -- cycle
  (0.0781,0.1353) -- (0.0938,0.1353) -- (0.0859,0.1488) -- cycle (0.0938,0.1353) -- (0.1094,0.1353) -- (0.1016,0.1488) -- cycle (0.0859,0.1488) -- (0.1016,0.1488) -- (0.0938,0.1624) -- cycle (0.125,0.1083) -- (0.1406,0.1083) -- (0.1328,0.1218) -- cycle
  (0.1406,0.1083) -- (0.1562,0.1083) -- (0.1484,0.1218) -- cycle (0.1328,0.1218) -- (0.1484,0.1218) -- (0.1406,0.1353) -- cycle (0.1562,0.1083) -- (0.1719,0.1083) -- (0.1641,0.1218) -- cycle (0.1719,0.1083) -- (0.1875,0.1083) -- (0.1797,0.1218) -- cycle
  (0.1641,0.1218) -- (0.1797,0.1218) -- (0.1719,0.1353) -- cycle (0.1406,0.1353) -- (0.1562,0.1353) -- (0.1484,0.1488) -- cycle (0.1562,0.1353) -- (0.1719,0.1353) -- (0.1641,0.1488) -- cycle (0.1484,0.1488) -- (0.1641,0.1488) -- (0.1562,0.1624) -- cycle
  (0.0938,0.1624) -- (0.1094,0.1624) -- (0.1016,0.1759) -- cycle (0.1094,0.1624) -- (0.125,0.1624) -- (0.1172,0.1759) -- cycle (0.1016,0.1759) -- (0.1172,0.1759) -- (0.1094,0.1894) -- cycle (0.125,0.1624) -- (0.1406,0.1624) -- (0.1328,0.1759) -- cycle
  (0.1406,0.1624) -- (0.1562,0.1624) -- (0.1484,0.1759) -- cycle (0.1328,0.1759) -- (0.1484,0.1759) -- (0.1406,0.1894) -- cycle (0.1094,0.1894) -- (0.125,0.1894) -- (0.1172,0.203) -- cycle (0.125,0.1894) -- (0.1406,0.1894) -- (0.1328,0.203) -- cycle
  (0.1172,0.203) -- (0.1328,0.203) -- (0.125,0.2165) -- cycle (0.25,0) -- (0.2656,0) -- (0.2578,0.0135) -- cycle (0.2656,0) -- (0.2812,0) -- (0.2734,0.0135) -- cycle (0.2578,0.0135) -- (0.2734,0.0135) -- (0.2656,0.0271) -- cycle
  (0.2812,0) -- (0.2969,0) -- (0.2891,0.0135) -- cycle (0.2969,0) -- (0.3125,0) -- (0.3047,0.0135) -- cycle (0.2891,0.0135) -- (0.3047,0.0135) -- (0.2969,0.0271) -- cycle (0.2656,0.0271) -- (0.2812,0.0271) -- (0.2734,0.0406) -- cycle
  (0.2812,0.0271) -- (0.2969,0.0271) -- (0.2891,0.0406) -- cycle (0.2734,0.0406) -- (0.2891,0.0406) -- (0.2812,0.0541) -- cycle (0.3125,0) -- (0.3281,0) -- (0.3203,0.0135) -- cycle (0.3281,0) -- (0.3438,0) -- (0.3359,0.0135) -- cycle
  (0.3203,0.0135) -- (0.3359,0.0135) -- (0.3281,0.0271) -- cycle (0.3438,0) -- (0.3594,0) -- (0.3516,0.0135) -- cycle (0.3594,0) -- (0.375,0) -- (0.3672,0.0135) -- cycle (0.3516,0.0135) -- (0.3672,0.0135) -- (0.3594,0.0271) -- cycle
  (0.3281,0.0271) -- (0.3438,0.0271) -- (0.3359,0.0406) -- cycle (0.3438,0.0271) -- (0.3594,0.0271) -- (0.3516,0.0406) -- cycle (0.3359,0.0406) -- (0.3516,0.0406) -- (0.3438,0.0541) -- cycle (0.2812,0.0541) -- (0.2969,0.0541) -- (0.2891,0.0677) -- cycle
  (0.2969,0.0541) -- (0.3125,0.0541) -- (0.3047,0.0677) -- cycle (0.2891,0.0677) -- (0.3047,0.0677) -- (0.2969,0.0812) -- cycle (0.3125,0.0541) -- (0.3281,0.0541) -- (0.3203,0.0677) -- cycle (0.3281,0.0541) -- (0.3438,0.0541) -- (0.3359,0.0677) -- cycle
  (0.3203,0.0677) -- (0.3359,0.0677) -- (0.3281,0.0812) -- cycle (0.2969,0.0812) -- (0.3125,0.0812) -- (0.3047,0.0947) -- cycle (0.3125,0.0812) -- (0.3281,0.0812) -- (0.3203,0.0947) -- cycle (0.3047,0.0947) -- (0.3203,0.0947) -- (0.3125,0.1083) -- cycle
  (0.375,0) -- (0.3906,0) -- (0.3828,0.0135) -- cycle (0.3906,0) -- (0.4062,0) -- (0.3984,0.0135) -- cycle (0.3828,0.0135) -- (0.3984,0.0135) -- (0.3906,0.0271) -- cycle (0.4062,0) -- (0.4219,0) -- (0.4141,0.0135) -- cycle
  (0.4219,0) -- (0.4375,0) -- (0.4297,0.0135) -- cycle (0.4141,0.0135) -- (0.4297,0.0135) -- (0.4219,0.0271) -- cycle (0.3906,0.0271) -- (0.4062,0.0271) -- (0.3984,0.0406) -- cycle (0.4062,0.0271) -- (0.4219,0.0271) -- (0.4141,0.0406) -- cycle
  (0.3984,0.0406) -- (0.4141,0.0406) -- (0.4062,0.0541) -- cycle (0.4375,0) -- (0.4531,0) -- (0.4453,0.0135) -- cycle (0.4531,0) -- (0.4688,0) -- (0.4609,0.0135) -- cycle (0.4453,0.0135) -- (0.4609,0.0135) -- (0.4531,0.0271) -- cycle
  (0.4688,0) -- (0.4844,0) -- (0.4766,0.0135) -- cycle (0.4844,0) -- (0.5,0) -- (0.4922,0.0135) -- cycle (0.4766,0.0135) -- (0.4922,0.0135) -- (0.4844,0.0271) -- cycle (0.4531,0.0271) -- (0.4688,0.0271) -- (0.4609,0.0406) -- cycle
  (0.4688,0.0271) -- (0.4844,0.0271) -- (0.4766,0.0406) -- cycle (0.4609,0.0406) -- (0.4766,0.0406) -- (0.4688,0.0541) -- cycle (0.4062,0.0541) -- (0.4219,0.0541) -- (0.4141,0.0677) -- cycle (0.4219,0.0541) -- (0.4375,0.0541) -- (0.4297,0.0677) -- cycle
  (0.4141,0.0677) -- (0.4297,0.0677) -- (0.4219,0.0812) -- cycle (0.4375,0.0541) -- (0.4531,0.0541) -- (0.4453,0.0677) -- cycle (0.4531,0.0541) -- (0.4688,0.0541) -- (0.4609,0.0677) -- cycle (0.4453,0.0677) -- (0.4609,0.0677) -- (0.4531,0.0812) -- cycle
  (0.4219,0.0812) -- (0.4375,0.0812) -- (0.4297,0.0947) -- cycle (0.4375,0.0812) -- (0.4531,0.0812) -- (0.4453,0.0947) -- cycle (0.4297,0.0947) -- (0.4453,0.0947) -- (0.4375,0.1083) -- cycle (0.3125,0.1083) -- (0.3281,0.1083) -- (0.3203,0.1218) -- cycle
  (0.3281,0.1083) -- (0.3438,0.1083) -- (0.3359,0.1218) -- cycle (0.3203,0.1218) -- (0.3359,0.1218) -- (0.3281,0.1353) -- cycle (0.3438,0.1083) -- (0.3594,0.1083) -- (0.3516,0.1218) -- cycle (0.3594,0.1083) -- (0.375,0.1083) -- (0.3672,0.1218) -- cycle
  (0.3516,0.1218) -- (0.3672,0.1218) -- (0.3594,0.1353) -- cycle (0.3281,0.1353) -- (0.3438,0.1353) -- (0.3359,0.1488) -- cycle (0.3438,0.1353) -- (0.3594,0.1353) -- (0.3516,0.1488) -- cycle (0.3359,0.1488) -- (0.3516,0.1488) -- (0.3438,0.1624) -- cycle
  (0.375,0.1083) -- (0.3906,0.1083) -- (0.3828,0.1218) -- cycle (0.3906,0.1083) -- (0.4062,0.1083) -- (0.3984,0.1218) -- cycle (0.3828,0.1218) -- (0.3984,0.1218) -- (0.3906,0.1353) -- cycle (0.4062,0.1083) -- (0.4219,0.1083) -- (0.4141,0.1218) -- cycle
  (0.4219,0.1083) -- (0.4375,0.1083) -- (0.4297,0.1218) -- cycle (0.4141,0.1218) -- (0.4297,0.1218) -- (0.4219,0.1353) -- cycle (0.3906,0.1353) -- (0.4062,0.1353) -- (0.3984,0.1488) -- cycle (0.4062,0.1353) -- (0.4219,0.1353) -- (0.4141,0.1488) -- cycle
  (0.3984,0.1488) -- (0.4141,0.1488) -- (0.4062,0.1624) -- cycle (0.3438,0.1624) -- (0.3594,0.1624) -- (0.3516,0.1759) -- cycle (0.3594,0.1624) -- (0.375,0.1624) -- (0.3672,0.1759) -- cycle (0.3516,0.1759) -- (0.3672,0.1759) -- (0.3594,0.1894) -- cycle
  (0.375,0.1624) -- (0.3906,0.1624) -- (0.3828,0.1759) -- cycle (0.3906,0.1624) -- (0.4062,0.1624) -- (0.3984,0.1759) -- cycle (0.3828,0.1759) -- (0.3984,0.1759) -- (0.3906,0.1894) -- cycle (0.3594,0.1894) -- (0.375,0.1894) -- (0.3672,0.203) -- cycle
  (0.375,0.1894) -- (0.3906,0.1894) -- (0.3828,0.203) -- cycle (0.3672,0.203) -- (0.3828,0.203) -- (0.375,0.2165) -- cycle (0.125,0.2165) -- (0.1406,0.2165) -- (0.1328,0.23) -- cycle (0.1406,0.2165) -- (0.1562,0.2165) -- (0.1484,0.23) -- cycle
  (0.1328,0.23) -- (0.1484,0.23) -- (0.1406,0.2436) -- cycle (0.1562,0.2165) -- (0.1719,0.2165) -- (0.1641,0.23) -- cycle (0.1719,0.2165) -- (0.1875,0.2165) -- (0.1797,0.23) -- cycle (0.1641,0.23) -- (0.1797,0.23) -- (0.1719,0.2436) -- cycle
  (0.1406,0.2436) -- (0.1562,0.2436) -- (0.1484,0.2571) -- cycle (0.1562,0.2436) -- (0.1719,0.2436) -- (0.1641,0.2571) -- cycle (0.1484,0.2571) -- (0.1641,0.2571) -- (0.1562,0.2706) -- cycle (0.1875,0.2165) -- (0.2031,0.2165) -- (0.1953,0.23) -- cycle
  (0.2031,0.2165) -- (0.2188,0.2165) -- (0.2109,0.23) -- cycle (0.1953,0.23) -- (0.2109,0.23) -- (0.2031,0.2436) -- cycle (0.2188,0.2165) -- (0.2344,0.2165) -- (0.2266,0.23) -- cycle (0.2344,0.2165) -- (0.25,0.2165) -- (0.2422,0.23) -- cycle
  (0.2266,0.23) -- (0.2422,0.23) -- (0.2344,0.2436) -- cycle (0.2031,0.2436) -- (0.2188,0.2436) -- (0.2109,0.2571) -- cycle (0.2188,0.2436) -- (0.2344,0.2436) -- (0.2266,0.2571) -- cycle (0.2109,0.2571) -- (0.2266,0.2571) -- (0.2188,0.2706) -- cycle
  (0.1562,0.2706) -- (0.1719,0.2706) -- (0.1641,0.2842) -- cycle (0.1719,0.2706) -- (0.1875,0.2706) -- (0.1797,0.2842) -- cycle (0.1641,0.2842) -- (0.1797,0.2842) -- (0.1719,0.2977) -- cycle (0.1875,0.2706) -- (0.2031,0.2706) -- (0.1953,0.2842) -- cycle
  (0.2031,0.2706) -- (0.2188,0.2706) -- (0.2109,0.2842) -- cycle (0.1953,0.2842) -- (0.2109,0.2842) -- (0.2031,0.2977) -- cycle (0.1719,0.2977) -- (0.1875,0.2977) -- (0.1797,0.3112) -- cycle (0.1875,0.2977) -- (0.2031,0.2977) -- (0.1953,0.3112) -- cycle
  (0.1797,0.3112) -- (0.1953,0.3112) -- (0.1875,0.3248) -- cycle (0.25,0.2165) -- (0.2656,0.2165) -- (0.2578,0.23) -- cycle (0.2656,0.2165) -- (0.2812,0.2165) -- (0.2734,0.23) -- cycle (0.2578,0.23) -- (0.2734,0.23) -- (0.2656,0.2436) -- cycle
  (0.2812,0.2165) -- (0.2969,0.2165) -- (0.2891,0.23) -- cycle (0.2969,0.2165) -- (0.3125,0.2165) -- (0.3047,0.23) -- cycle (0.2891,0.23) -- (0.3047,0.23) -- (0.2969,0.2436) -- cycle (0.2656,0.2436) -- (0.2812,0.2436) -- (0.2734,0.2571) -- cycle
  (0.2812,0.2436) -- (0.2969,0.2436) -- (0.2891,0.2571) -- cycle (0.2734,0.2571) -- (0.2891,0.2571) -- (0.2812,0.2706) -- cycle (0.3125,0.2165) -- (0.3281,0.2165) -- (0.3203,0.23) -- cycle (0.3281,0.2165) -- (0.3438,0.2165) -- (0.3359,0.23) -- cycle
  (0.3203,0.23) -- (0.3359,0.23) -- (0.3281,0.2436) -- cycle (0.3438,0.2165) -- (0.3594,0.2165) -- (0.3516,0.23) -- cycle (0.3594,0.2165) -- (0.375,0.2165) -- (0.3672,0.23) -- cycle (0.3516,0.23) -- (0.3672,0.23) -- (0.3594,0.2436) -- cycle
  (0.3281,0.2436) -- (0.3438,0.2436) -- (0.3359,0.2571) -- cycle (0.3438,0.2436) -- (0.3594,0.2436) -- (0.3516,0.2571) -- cycle (0.3359,0.2571) -- (0.3516,0.2571) -- (0.3438,0.2706) -- cycle (0.2812,0.2706) -- (0.2969,0.2706) -- (0.2891,0.2842) -- cycle
  (0.2969,0.2706) -- (0.3125,0.2706) -- (0.3047,0.2842) -- cycle (0.2891,0.2842) -- (0.3047,0.2842) -- (0.2969,0.2977) -- cycle (0.3125,0.2706) -- (0.3281,0.2706) -- (0.3203,0.2842) -- cycle (0.3281,0.2706) -- (0.3438,0.2706) -- (0.3359,0.2842) -- cycle
  (0.3203,0.2842) -- (0.3359,0.2842) -- (0.3281,0.2977) -- cycle (0.2969,0.2977) -- (0.3125,0.2977) -- (0.3047,0.3112) -- cycle (0.3125,0.2977) -- (0.3281,0.2977) -- (0.3203,0.3112) -- cycle (0.3047,0.3112) -- (0.3203,0.3112) -- (0.3125,0.3248) -- cycle
  (0.1875,0.3248) -- (0.2031,0.3248) -- (0.1953,0.3383) -- cycle (0.2031,0.3248) -- (0.2188,0.3248) -- (0.2109,0.3383) -- cycle (0.1953,0.3383) -- (0.2109,0.3383) -- (0.2031,0.3518) -- cycle (0.2188,0.3248) -- (0.2344,0.3248) -- (0.2266,0.3383) -- cycle
  (0.2344,0.3248) -- (0.25,0.3248) -- (0.2422,0.3383) -- cycle (0.2266,0.3383) -- (0.2422,0.3383) -- (0.2344,0.3518) -- cycle (0.2031,0.3518) -- (0.2188,0.3518) -- (0.2109,0.3654) -- cycle (0.2188,0.3518) -- (0.2344,0.3518) -- (0.2266,0.3654) -- cycle
  (0.2109,0.3654) -- (0.2266,0.3654) -- (0.2188,0.3789) -- cycle (0.25,0.3248) -- (0.2656,0.3248) -- (0.2578,0.3383) -- cycle (0.2656,0.3248) -- (0.2812,0.3248) -- (0.2734,0.3383) -- cycle (0.2578,0.3383) -- (0.2734,0.3383) -- (0.2656,0.3518) -- cycle
  (0.2812,0.3248) -- (0.2969,0.3248) -- (0.2891,0.3383) -- cycle (0.2969,0.3248) -- (0.3125,0.3248) -- (0.3047,0.3383) -- cycle (0.2891,0.3383) -- (0.3047,0.3383) -- (0.2969,0.3518) -- cycle (0.2656,0.3518) -- (0.2812,0.3518) -- (0.2734,0.3654) -- cycle
  (0.2812,0.3518) -- (0.2969,0.3518) -- (0.2891,0.3654) -- cycle (0.2734,0.3654) -- (0.2891,0.3654) -- (0.2812,0.3789) -- cycle (0.2188,0.3789) -- (0.2344,0.3789) -- (0.2266,0.3924) -- cycle (0.2344,0.3789) -- (0.25,0.3789) -- (0.2422,0.3924) -- cycle
  (0.2266,0.3924) -- (0.2422,0.3924) -- (0.2344,0.4059) -- cycle (0.25,0.3789) -- (0.2656,0.3789) -- (0.2578,0.3924) -- cycle (0.2656,0.3789) -- (0.2812,0.3789) -- (0.2734,0.3924) -- cycle (0.2578,0.3924) -- (0.2734,0.3924) -- (0.2656,0.4059) -- cycle
  (0.2344,0.4059) -- (0.25,0.4059) -- (0.2422,0.4195) -- cycle (0.25,0.4059) -- (0.2656,0.4059) -- (0.2578,0.4195) -- cycle (0.2422,0.4195) -- (0.2578,0.4195) -- (0.25,0.433) -- cycle;
\fill[piece2]
  (0.5,0) -- (0.5156,0) -- (0.5078,0.0135) -- cycle (0.5156,0) -- (0.5312,0) -- (0.5234,0.0135) -- cycle (0.5078,0.0135) -- (0.5234,0.0135) -- (0.5156,0.0271) -- cycle (0.5312,0) -- (0.5469,0) -- (0.5391,0.0135) -- cycle
  (0.5469,0) -- (0.5625,0) -- (0.5547,0.0135) -- cycle (0.5391,0.0135) -- (0.5547,0.0135) -- (0.5469,0.0271) -- cycle (0.5156,0.0271) -- (0.5312,0.0271) -- (0.5234,0.0406) -- cycle (0.5312,0.0271) -- (0.5469,0.0271) -- (0.5391,0.0406) -- cycle
  (0.5234,0.0406) -- (0.5391,0.0406) -- (0.5312,0.0541) -- cycle (0.5625,0) -- (0.5781,0) -- (0.5703,0.0135) -- cycle (0.5781,0) -- (0.5938,0) -- (0.5859,0.0135) -- cycle (0.5703,0.0135) -- (0.5859,0.0135) -- (0.5781,0.0271) -- cycle
  (0.5938,0) -- (0.6094,0) -- (0.6016,0.0135) -- cycle (0.6094,0) -- (0.625,0) -- (0.6172,0.0135) -- cycle (0.6016,0.0135) -- (0.6172,0.0135) -- (0.6094,0.0271) -- cycle (0.5781,0.0271) -- (0.5938,0.0271) -- (0.5859,0.0406) -- cycle
  (0.5938,0.0271) -- (0.6094,0.0271) -- (0.6016,0.0406) -- cycle (0.5859,0.0406) -- (0.6016,0.0406) -- (0.5938,0.0541) -- cycle (0.5312,0.0541) -- (0.5469,0.0541) -- (0.5391,0.0677) -- cycle (0.5469,0.0541) -- (0.5625,0.0541) -- (0.5547,0.0677) -- cycle
  (0.5391,0.0677) -- (0.5547,0.0677) -- (0.5469,0.0812) -- cycle (0.5625,0.0541) -- (0.5781,0.0541) -- (0.5703,0.0677) -- cycle (0.5781,0.0541) -- (0.5938,0.0541) -- (0.5859,0.0677) -- cycle (0.5703,0.0677) -- (0.5859,0.0677) -- (0.5781,0.0812) -- cycle
  (0.5469,0.0812) -- (0.5625,0.0812) -- (0.5547,0.0947) -- cycle (0.5625,0.0812) -- (0.5781,0.0812) -- (0.5703,0.0947) -- cycle (0.5547,0.0947) -- (0.5703,0.0947) -- (0.5625,0.1083) -- cycle (0.625,0) -- (0.6406,0) -- (0.6328,0.0135) -- cycle
  (0.6406,0) -- (0.6562,0) -- (0.6484,0.0135) -- cycle (0.6328,0.0135) -- (0.6484,0.0135) -- (0.6406,0.0271) -- cycle (0.6562,0) -- (0.6719,0) -- (0.6641,0.0135) -- cycle (0.6719,0) -- (0.6875,0) -- (0.6797,0.0135) -- cycle
  (0.6641,0.0135) -- (0.6797,0.0135) -- (0.6719,0.0271) -- cycle (0.6406,0.0271) -- (0.6562,0.0271) -- (0.6484,0.0406) -- cycle (0.6562,0.0271) -- (0.6719,0.0271) -- (0.6641,0.0406) -- cycle (0.6484,0.0406) -- (0.6641,0.0406) -- (0.6562,0.0541) -- cycle
  (0.6875,0) -- (0.7031,0) -- (0.6953,0.0135) -- cycle (0.7031,0) -- (0.7188,0) -- (0.7109,0.0135) -- cycle (0.6953,0.0135) -- (0.7109,0.0135) -- (0.7031,0.0271) -- cycle (0.7188,0) -- (0.7344,0) -- (0.7266,0.0135) -- cycle
  (0.7344,0) -- (0.75,0) -- (0.7422,0.0135) -- cycle (0.7266,0.0135) -- (0.7422,0.0135) -- (0.7344,0.0271) -- cycle (0.7031,0.0271) -- (0.7188,0.0271) -- (0.7109,0.0406) -- cycle (0.7188,0.0271) -- (0.7344,0.0271) -- (0.7266,0.0406) -- cycle
  (0.7109,0.0406) -- (0.7266,0.0406) -- (0.7188,0.0541) -- cycle (0.6562,0.0541) -- (0.6719,0.0541) -- (0.6641,0.0677) -- cycle (0.6719,0.0541) -- (0.6875,0.0541) -- (0.6797,0.0677) -- cycle (0.6641,0.0677) -- (0.6797,0.0677) -- (0.6719,0.0812) -- cycle
  (0.6875,0.0541) -- (0.7031,0.0541) -- (0.6953,0.0677) -- cycle (0.7031,0.0541) -- (0.7188,0.0541) -- (0.7109,0.0677) -- cycle (0.6953,0.0677) -- (0.7109,0.0677) -- (0.7031,0.0812) -- cycle (0.6719,0.0812) -- (0.6875,0.0812) -- (0.6797,0.0947) -- cycle
  (0.6875,0.0812) -- (0.7031,0.0812) -- (0.6953,0.0947) -- cycle (0.6797,0.0947) -- (0.6953,0.0947) -- (0.6875,0.1083) -- cycle (0.5625,0.1083) -- (0.5781,0.1083) -- (0.5703,0.1218) -- cycle (0.5781,0.1083) -- (0.5938,0.1083) -- (0.5859,0.1218) -- cycle
  (0.5703,0.1218) -- (0.5859,0.1218) -- (0.5781,0.1353) -- cycle (0.5938,0.1083) -- (0.6094,0.1083) -- (0.6016,0.1218) -- cycle (0.6094,0.1083) -- (0.625,0.1083) -- (0.6172,0.1218) -- cycle (0.6016,0.1218) -- (0.6172,0.1218) -- (0.6094,0.1353) -- cycle
  (0.5781,0.1353) -- (0.5938,0.1353) -- (0.5859,0.1488) -- cycle (0.5938,0.1353) -- (0.6094,0.1353) -- (0.6016,0.1488) -- cycle (0.5859,0.1488) -- (0.6016,0.1488) -- (0.5938,0.1624) -- cycle (0.625,0.1083) -- (0.6406,0.1083) -- (0.6328,0.1218) -- cycle
  (0.6406,0.1083) -- (0.6562,0.1083) -- (0.6484,0.1218) -- cycle (0.6328,0.1218) -- (0.6484,0.1218) -- (0.6406,0.1353) -- cycle (0.6562,0.1083) -- (0.6719,0.1083) -- (0.6641,0.1218) -- cycle (0.6719,0.1083) -- (0.6875,0.1083) -- (0.6797,0.1218) -- cycle
  (0.6641,0.1218) -- (0.6797,0.1218) -- (0.6719,0.1353) -- cycle (0.6406,0.1353) -- (0.6562,0.1353) -- (0.6484,0.1488) -- cycle (0.6562,0.1353) -- (0.6719,0.1353) -- (0.6641,0.1488) -- cycle (0.6484,0.1488) -- (0.6641,0.1488) -- (0.6562,0.1624) -- cycle
  (0.5938,0.1624) -- (0.6094,0.1624) -- (0.6016,0.1759) -- cycle (0.6094,0.1624) -- (0.625,0.1624) -- (0.6172,0.1759) -- cycle (0.6016,0.1759) -- (0.6172,0.1759) -- (0.6094,0.1894) -- cycle (0.625,0.1624) -- (0.6406,0.1624) -- (0.6328,0.1759) -- cycle
  (0.6406,0.1624) -- (0.6562,0.1624) -- (0.6484,0.1759) -- cycle (0.6328,0.1759) -- (0.6484,0.1759) -- (0.6406,0.1894) -- cycle (0.6094,0.1894) -- (0.625,0.1894) -- (0.6172,0.203) -- cycle (0.625,0.1894) -- (0.6406,0.1894) -- (0.6328,0.203) -- cycle
  (0.6172,0.203) -- (0.6328,0.203) -- (0.625,0.2165) -- cycle (0.75,0) -- (0.7656,0) -- (0.7578,0.0135) -- cycle (0.7656,0) -- (0.7812,0) -- (0.7734,0.0135) -- cycle (0.7578,0.0135) -- (0.7734,0.0135) -- (0.7656,0.0271) -- cycle
  (0.7812,0) -- (0.7969,0) -- (0.7891,0.0135) -- cycle (0.7969,0) -- (0.8125,0) -- (0.8047,0.0135) -- cycle (0.7891,0.0135) -- (0.8047,0.0135) -- (0.7969,0.0271) -- cycle (0.7656,0.0271) -- (0.7812,0.0271) -- (0.7734,0.0406) -- cycle
  (0.7812,0.0271) -- (0.7969,0.0271) -- (0.7891,0.0406) -- cycle (0.7734,0.0406) -- (0.7891,0.0406) -- (0.7812,0.0541) -- cycle (0.8125,0) -- (0.8281,0) -- (0.8203,0.0135) -- cycle (0.8281,0) -- (0.8438,0) -- (0.8359,0.0135) -- cycle
  (0.8203,0.0135) -- (0.8359,0.0135) -- (0.8281,0.0271) -- cycle (0.8438,0) -- (0.8594,0) -- (0.8516,0.0135) -- cycle (0.8594,0) -- (0.875,0) -- (0.8672,0.0135) -- cycle (0.8516,0.0135) -- (0.8672,0.0135) -- (0.8594,0.0271) -- cycle
  (0.8281,0.0271) -- (0.8438,0.0271) -- (0.8359,0.0406) -- cycle (0.8438,0.0271) -- (0.8594,0.0271) -- (0.8516,0.0406) -- cycle (0.8359,0.0406) -- (0.8516,0.0406) -- (0.8438,0.0541) -- cycle (0.7812,0.0541) -- (0.7969,0.0541) -- (0.7891,0.0677) -- cycle
  (0.7969,0.0541) -- (0.8125,0.0541) -- (0.8047,0.0677) -- cycle (0.7891,0.0677) -- (0.8047,0.0677) -- (0.7969,0.0812) -- cycle (0.8125,0.0541) -- (0.8281,0.0541) -- (0.8203,0.0677) -- cycle (0.8281,0.0541) -- (0.8438,0.0541) -- (0.8359,0.0677) -- cycle
  (0.8203,0.0677) -- (0.8359,0.0677) -- (0.8281,0.0812) -- cycle (0.7969,0.0812) -- (0.8125,0.0812) -- (0.8047,0.0947) -- cycle (0.8125,0.0812) -- (0.8281,0.0812) -- (0.8203,0.0947) -- cycle (0.8047,0.0947) -- (0.8203,0.0947) -- (0.8125,0.1083) -- cycle
  (0.875,0) -- (0.8906,0) -- (0.8828,0.0135) -- cycle (0.8906,0) -- (0.9062,0) -- (0.8984,0.0135) -- cycle (0.8828,0.0135) -- (0.8984,0.0135) -- (0.8906,0.0271) -- cycle (0.9062,0) -- (0.9219,0) -- (0.9141,0.0135) -- cycle
  (0.9219,0) -- (0.9375,0) -- (0.9297,0.0135) -- cycle (0.9141,0.0135) -- (0.9297,0.0135) -- (0.9219,0.0271) -- cycle (0.8906,0.0271) -- (0.9062,0.0271) -- (0.8984,0.0406) -- cycle (0.9062,0.0271) -- (0.9219,0.0271) -- (0.9141,0.0406) -- cycle
  (0.8984,0.0406) -- (0.9141,0.0406) -- (0.9062,0.0541) -- cycle (0.9375,0) -- (0.9531,0) -- (0.9453,0.0135) -- cycle (0.9531,0) -- (0.9688,0) -- (0.9609,0.0135) -- cycle (0.9453,0.0135) -- (0.9609,0.0135) -- (0.9531,0.0271) -- cycle
  (0.9688,0) -- (0.9844,0) -- (0.9766,0.0135) -- cycle (0.9844,0) -- (1,0) -- (0.9922,0.0135) -- cycle (0.9766,0.0135) -- (0.9922,0.0135) -- (0.9844,0.0271) -- cycle (0.9531,0.0271) -- (0.9688,0.0271) -- (0.9609,0.0406) -- cycle
  (0.9688,0.0271) -- (0.9844,0.0271) -- (0.9766,0.0406) -- cycle (0.9609,0.0406) -- (0.9766,0.0406) -- (0.9688,0.0541) -- cycle (0.9062,0.0541) -- (0.9219,0.0541) -- (0.9141,0.0677) -- cycle (0.9219,0.0541) -- (0.9375,0.0541) -- (0.9297,0.0677) -- cycle
  (0.9141,0.0677) -- (0.9297,0.0677) -- (0.9219,0.0812) -- cycle (0.9375,0.0541) -- (0.9531,0.0541) -- (0.9453,0.0677) -- cycle (0.9531,0.0541) -- (0.9688,0.0541) -- (0.9609,0.0677) -- cycle (0.9453,0.0677) -- (0.9609,0.0677) -- (0.9531,0.0812) -- cycle
  (0.9219,0.0812) -- (0.9375,0.0812) -- (0.9297,0.0947) -- cycle (0.9375,0.0812) -- (0.9531,0.0812) -- (0.9453,0.0947) -- cycle (0.9297,0.0947) -- (0.9453,0.0947) -- (0.9375,0.1083) -- cycle (0.8125,0.1083) -- (0.8281,0.1083) -- (0.8203,0.1218) -- cycle
  (0.8281,0.1083) -- (0.8438,0.1083) -- (0.8359,0.1218) -- cycle (0.8203,0.1218) -- (0.8359,0.1218) -- (0.8281,0.1353) -- cycle (0.8438,0.1083) -- (0.8594,0.1083) -- (0.8516,0.1218) -- cycle (0.8594,0.1083) -- (0.875,0.1083) -- (0.8672,0.1218) -- cycle
  (0.8516,0.1218) -- (0.8672,0.1218) -- (0.8594,0.1353) -- cycle (0.8281,0.1353) -- (0.8438,0.1353) -- (0.8359,0.1488) -- cycle (0.8438,0.1353) -- (0.8594,0.1353) -- (0.8516,0.1488) -- cycle (0.8359,0.1488) -- (0.8516,0.1488) -- (0.8438,0.1624) -- cycle
  (0.875,0.1083) -- (0.8906,0.1083) -- (0.8828,0.1218) -- cycle (0.8906,0.1083) -- (0.9062,0.1083) -- (0.8984,0.1218) -- cycle (0.8828,0.1218) -- (0.8984,0.1218) -- (0.8906,0.1353) -- cycle (0.9062,0.1083) -- (0.9219,0.1083) -- (0.9141,0.1218) -- cycle
  (0.9219,0.1083) -- (0.9375,0.1083) -- (0.9297,0.1218) -- cycle (0.9141,0.1218) -- (0.9297,0.1218) -- (0.9219,0.1353) -- cycle (0.8906,0.1353) -- (0.9062,0.1353) -- (0.8984,0.1488) -- cycle (0.9062,0.1353) -- (0.9219,0.1353) -- (0.9141,0.1488) -- cycle
  (0.8984,0.1488) -- (0.9141,0.1488) -- (0.9062,0.1624) -- cycle (0.8438,0.1624) -- (0.8594,0.1624) -- (0.8516,0.1759) -- cycle (0.8594,0.1624) -- (0.875,0.1624) -- (0.8672,0.1759) -- cycle (0.8516,0.1759) -- (0.8672,0.1759) -- (0.8594,0.1894) -- cycle
  (0.875,0.1624) -- (0.8906,0.1624) -- (0.8828,0.1759) -- cycle (0.8906,0.1624) -- (0.9062,0.1624) -- (0.8984,0.1759) -- cycle (0.8828,0.1759) -- (0.8984,0.1759) -- (0.8906,0.1894) -- cycle (0.8594,0.1894) -- (0.875,0.1894) -- (0.8672,0.203) -- cycle
  (0.875,0.1894) -- (0.8906,0.1894) -- (0.8828,0.203) -- cycle (0.8672,0.203) -- (0.8828,0.203) -- (0.875,0.2165) -- cycle (0.625,0.2165) -- (0.6406,0.2165) -- (0.6328,0.23) -- cycle (0.6406,0.2165) -- (0.6562,0.2165) -- (0.6484,0.23) -- cycle
  (0.6328,0.23) -- (0.6484,0.23) -- (0.6406,0.2436) -- cycle (0.6562,0.2165) -- (0.6719,0.2165) -- (0.6641,0.23) -- cycle (0.6719,0.2165) -- (0.6875,0.2165) -- (0.6797,0.23) -- cycle (0.6641,0.23) -- (0.6797,0.23) -- (0.6719,0.2436) -- cycle
  (0.6406,0.2436) -- (0.6562,0.2436) -- (0.6484,0.2571) -- cycle (0.6562,0.2436) -- (0.6719,0.2436) -- (0.6641,0.2571) -- cycle (0.6484,0.2571) -- (0.6641,0.2571) -- (0.6562,0.2706) -- cycle (0.6875,0.2165) -- (0.7031,0.2165) -- (0.6953,0.23) -- cycle
  (0.7031,0.2165) -- (0.7188,0.2165) -- (0.7109,0.23) -- cycle (0.6953,0.23) -- (0.7109,0.23) -- (0.7031,0.2436) -- cycle (0.7188,0.2165) -- (0.7344,0.2165) -- (0.7266,0.23) -- cycle (0.7344,0.2165) -- (0.75,0.2165) -- (0.7422,0.23) -- cycle
  (0.7266,0.23) -- (0.7422,0.23) -- (0.7344,0.2436) -- cycle (0.7031,0.2436) -- (0.7188,0.2436) -- (0.7109,0.2571) -- cycle (0.7188,0.2436) -- (0.7344,0.2436) -- (0.7266,0.2571) -- cycle (0.7109,0.2571) -- (0.7266,0.2571) -- (0.7188,0.2706) -- cycle
  (0.6562,0.2706) -- (0.6719,0.2706) -- (0.6641,0.2842) -- cycle (0.6719,0.2706) -- (0.6875,0.2706) -- (0.6797,0.2842) -- cycle (0.6641,0.2842) -- (0.6797,0.2842) -- (0.6719,0.2977) -- cycle (0.6875,0.2706) -- (0.7031,0.2706) -- (0.6953,0.2842) -- cycle
  (0.7031,0.2706) -- (0.7188,0.2706) -- (0.7109,0.2842) -- cycle (0.6953,0.2842) -- (0.7109,0.2842) -- (0.7031,0.2977) -- cycle (0.6719,0.2977) -- (0.6875,0.2977) -- (0.6797,0.3112) -- cycle (0.6875,0.2977) -- (0.7031,0.2977) -- (0.6953,0.3112) -- cycle
  (0.6797,0.3112) -- (0.6953,0.3112) -- (0.6875,0.3248) -- cycle (0.75,0.2165) -- (0.7656,0.2165) -- (0.7578,0.23) -- cycle (0.7656,0.2165) -- (0.7812,0.2165) -- (0.7734,0.23) -- cycle (0.7578,0.23) -- (0.7734,0.23) -- (0.7656,0.2436) -- cycle
  (0.7812,0.2165) -- (0.7969,0.2165) -- (0.7891,0.23) -- cycle (0.7969,0.2165) -- (0.8125,0.2165) -- (0.8047,0.23) -- cycle (0.7891,0.23) -- (0.8047,0.23) -- (0.7969,0.2436) -- cycle (0.7656,0.2436) -- (0.7812,0.2436) -- (0.7734,0.2571) -- cycle
  (0.7812,0.2436) -- (0.7969,0.2436) -- (0.7891,0.2571) -- cycle (0.7734,0.2571) -- (0.7891,0.2571) -- (0.7812,0.2706) -- cycle (0.8125,0.2165) -- (0.8281,0.2165) -- (0.8203,0.23) -- cycle (0.8281,0.2165) -- (0.8438,0.2165) -- (0.8359,0.23) -- cycle
  (0.8203,0.23) -- (0.8359,0.23) -- (0.8281,0.2436) -- cycle (0.8438,0.2165) -- (0.8594,0.2165) -- (0.8516,0.23) -- cycle (0.8594,0.2165) -- (0.875,0.2165) -- (0.8672,0.23) -- cycle (0.8516,0.23) -- (0.8672,0.23) -- (0.8594,0.2436) -- cycle
  (0.8281,0.2436) -- (0.8438,0.2436) -- (0.8359,0.2571) -- cycle (0.8438,0.2436) -- (0.8594,0.2436) -- (0.8516,0.2571) -- cycle (0.8359,0.2571) -- (0.8516,0.2571) -- (0.8438,0.2706) -- cycle (0.7812,0.2706) -- (0.7969,0.2706) -- (0.7891,0.2842) -- cycle
  (0.7969,0.2706) -- (0.8125,0.2706) -- (0.8047,0.2842) -- cycle (0.7891,0.2842) -- (0.8047,0.2842) -- (0.7969,0.2977) -- cycle (0.8125,0.2706) -- (0.8281,0.2706) -- (0.8203,0.2842) -- cycle (0.8281,0.2706) -- (0.8438,0.2706) -- (0.8359,0.2842) -- cycle
  (0.8203,0.2842) -- (0.8359,0.2842) -- (0.8281,0.2977) -- cycle (0.7969,0.2977) -- (0.8125,0.2977) -- (0.8047,0.3112) -- cycle (0.8125,0.2977) -- (0.8281,0.2977) -- (0.8203,0.3112) -- cycle (0.8047,0.3112) -- (0.8203,0.3112) -- (0.8125,0.3248) -- cycle
  (0.6875,0.3248) -- (0.7031,0.3248) -- (0.6953,0.3383) -- cycle (0.7031,0.3248) -- (0.7188,0.3248) -- (0.7109,0.3383) -- cycle (0.6953,0.3383) -- (0.7109,0.3383) -- (0.7031,0.3518) -- cycle (0.7188,0.3248) -- (0.7344,0.3248) -- (0.7266,0.3383) -- cycle
  (0.7344,0.3248) -- (0.75,0.3248) -- (0.7422,0.3383) -- cycle (0.7266,0.3383) -- (0.7422,0.3383) -- (0.7344,0.3518) -- cycle (0.7031,0.3518) -- (0.7188,0.3518) -- (0.7109,0.3654) -- cycle (0.7188,0.3518) -- (0.7344,0.3518) -- (0.7266,0.3654) -- cycle
  (0.7109,0.3654) -- (0.7266,0.3654) -- (0.7188,0.3789) -- cycle (0.75,0.3248) -- (0.7656,0.3248) -- (0.7578,0.3383) -- cycle (0.7656,0.3248) -- (0.7812,0.3248) -- (0.7734,0.3383) -- cycle (0.7578,0.3383) -- (0.7734,0.3383) -- (0.7656,0.3518) -- cycle
  (0.7812,0.3248) -- (0.7969,0.3248) -- (0.7891,0.3383) -- cycle (0.7969,0.3248) -- (0.8125,0.3248) -- (0.8047,0.3383) -- cycle (0.7891,0.3383) -- (0.8047,0.3383) -- (0.7969,0.3518) -- cycle (0.7656,0.3518) -- (0.7812,0.3518) -- (0.7734,0.3654) -- cycle
  (0.7812,0.3518) -- (0.7969,0.3518) -- (0.7891,0.3654) -- cycle (0.7734,0.3654) -- (0.7891,0.3654) -- (0.7812,0.3789) -- cycle (0.7188,0.3789) -- (0.7344,0.3789) -- (0.7266,0.3924) -- cycle (0.7344,0.3789) -- (0.75,0.3789) -- (0.7422,0.3924) -- cycle
  (0.7266,0.3924) -- (0.7422,0.3924) -- (0.7344,0.4059) -- cycle (0.75,0.3789) -- (0.7656,0.3789) -- (0.7578,0.3924) -- cycle (0.7656,0.3789) -- (0.7812,0.3789) -- (0.7734,0.3924) -- cycle (0.7578,0.3924) -- (0.7734,0.3924) -- (0.7656,0.4059) -- cycle
  (0.7344,0.4059) -- (0.75,0.4059) -- (0.7422,0.4195) -- cycle (0.75,0.4059) -- (0.7656,0.4059) -- (0.7578,0.4195) -- cycle (0.7422,0.4195) -- (0.7578,0.4195) -- (0.75,0.433) -- cycle;
\fill[piece3]
  (0.25,0.433) -- (0.2656,0.433) -- (0.2578,0.4465) -- cycle (0.2656,0.433) -- (0.2812,0.433) -- (0.2734,0.4465) -- cycle (0.2578,0.4465) -- (0.2734,0.4465) -- (0.2656,0.4601) -- cycle (0.2812,0.433) -- (0.2969,0.433) -- (0.2891,0.4465) -- cycle
  (0.2969,0.433) -- (0.3125,0.433) -- (0.3047,0.4465) -- cycle (0.2891,0.4465) -- (0.3047,0.4465) -- (0.2969,0.4601) -- cycle (0.2656,0.4601) -- (0.2812,0.4601) -- (0.2734,0.4736) -- cycle (0.2812,0.4601) -- (0.2969,0.4601) -- (0.2891,0.4736) -- cycle
  (0.2734,0.4736) -- (0.2891,0.4736) -- (0.2812,0.4871) -- cycle (0.3125,0.433) -- (0.3281,0.433) -- (0.3203,0.4465) -- cycle (0.3281,0.433) -- (0.3438,0.433) -- (0.3359,0.4465) -- cycle (0.3203,0.4465) -- (0.3359,0.4465) -- (0.3281,0.4601) -- cycle
  (0.3438,0.433) -- (0.3594,0.433) -- (0.3516,0.4465) -- cycle (0.3594,0.433) -- (0.375,0.433) -- (0.3672,0.4465) -- cycle (0.3516,0.4465) -- (0.3672,0.4465) -- (0.3594,0.4601) -- cycle (0.3281,0.4601) -- (0.3438,0.4601) -- (0.3359,0.4736) -- cycle
  (0.3438,0.4601) -- (0.3594,0.4601) -- (0.3516,0.4736) -- cycle (0.3359,0.4736) -- (0.3516,0.4736) -- (0.3438,0.4871) -- cycle (0.2812,0.4871) -- (0.2969,0.4871) -- (0.2891,0.5007) -- cycle (0.2969,0.4871) -- (0.3125,0.4871) -- (0.3047,0.5007) -- cycle
  (0.2891,0.5007) -- (0.3047,0.5007) -- (0.2969,0.5142) -- cycle (0.3125,0.4871) -- (0.3281,0.4871) -- (0.3203,0.5007) -- cycle (0.3281,0.4871) -- (0.3438,0.4871) -- (0.3359,0.5007) -- cycle (0.3203,0.5007) -- (0.3359,0.5007) -- (0.3281,0.5142) -- cycle
  (0.2969,0.5142) -- (0.3125,0.5142) -- (0.3047,0.5277) -- cycle (0.3125,0.5142) -- (0.3281,0.5142) -- (0.3203,0.5277) -- cycle (0.3047,0.5277) -- (0.3203,0.5277) -- (0.3125,0.5413) -- cycle (0.375,0.433) -- (0.3906,0.433) -- (0.3828,0.4465) -- cycle
  (0.3906,0.433) -- (0.4062,0.433) -- (0.3984,0.4465) -- cycle (0.3828,0.4465) -- (0.3984,0.4465) -- (0.3906,0.4601) -- cycle (0.4062,0.433) -- (0.4219,0.433) -- (0.4141,0.4465) -- cycle (0.4219,0.433) -- (0.4375,0.433) -- (0.4297,0.4465) -- cycle
  (0.4141,0.4465) -- (0.4297,0.4465) -- (0.4219,0.4601) -- cycle (0.3906,0.4601) -- (0.4062,0.4601) -- (0.3984,0.4736) -- cycle (0.4062,0.4601) -- (0.4219,0.4601) -- (0.4141,0.4736) -- cycle (0.3984,0.4736) -- (0.4141,0.4736) -- (0.4062,0.4871) -- cycle
  (0.4375,0.433) -- (0.4531,0.433) -- (0.4453,0.4465) -- cycle (0.4531,0.433) -- (0.4688,0.433) -- (0.4609,0.4465) -- cycle (0.4453,0.4465) -- (0.4609,0.4465) -- (0.4531,0.4601) -- cycle (0.4688,0.433) -- (0.4844,0.433) -- (0.4766,0.4465) -- cycle
  (0.4844,0.433) -- (0.5,0.433) -- (0.4922,0.4465) -- cycle (0.4766,0.4465) -- (0.4922,0.4465) -- (0.4844,0.4601) -- cycle (0.4531,0.4601) -- (0.4688,0.4601) -- (0.4609,0.4736) -- cycle (0.4688,0.4601) -- (0.4844,0.4601) -- (0.4766,0.4736) -- cycle
  (0.4609,0.4736) -- (0.4766,0.4736) -- (0.4688,0.4871) -- cycle (0.4062,0.4871) -- (0.4219,0.4871) -- (0.4141,0.5007) -- cycle (0.4219,0.4871) -- (0.4375,0.4871) -- (0.4297,0.5007) -- cycle (0.4141,0.5007) -- (0.4297,0.5007) -- (0.4219,0.5142) -- cycle
  (0.4375,0.4871) -- (0.4531,0.4871) -- (0.4453,0.5007) -- cycle (0.4531,0.4871) -- (0.4688,0.4871) -- (0.4609,0.5007) -- cycle (0.4453,0.5007) -- (0.4609,0.5007) -- (0.4531,0.5142) -- cycle (0.4219,0.5142) -- (0.4375,0.5142) -- (0.4297,0.5277) -- cycle
  (0.4375,0.5142) -- (0.4531,0.5142) -- (0.4453,0.5277) -- cycle (0.4297,0.5277) -- (0.4453,0.5277) -- (0.4375,0.5413) -- cycle (0.3125,0.5413) -- (0.3281,0.5413) -- (0.3203,0.5548) -- cycle (0.3281,0.5413) -- (0.3438,0.5413) -- (0.3359,0.5548) -- cycle
  (0.3203,0.5548) -- (0.3359,0.5548) -- (0.3281,0.5683) -- cycle (0.3438,0.5413) -- (0.3594,0.5413) -- (0.3516,0.5548) -- cycle (0.3594,0.5413) -- (0.375,0.5413) -- (0.3672,0.5548) -- cycle (0.3516,0.5548) -- (0.3672,0.5548) -- (0.3594,0.5683) -- cycle
  (0.3281,0.5683) -- (0.3438,0.5683) -- (0.3359,0.5819) -- cycle (0.3438,0.5683) -- (0.3594,0.5683) -- (0.3516,0.5819) -- cycle (0.3359,0.5819) -- (0.3516,0.5819) -- (0.3438,0.5954) -- cycle (0.375,0.5413) -- (0.3906,0.5413) -- (0.3828,0.5548) -- cycle
  (0.3906,0.5413) -- (0.4062,0.5413) -- (0.3984,0.5548) -- cycle (0.3828,0.5548) -- (0.3984,0.5548) -- (0.3906,0.5683) -- cycle (0.4062,0.5413) -- (0.4219,0.5413) -- (0.4141,0.5548) -- cycle (0.4219,0.5413) -- (0.4375,0.5413) -- (0.4297,0.5548) -- cycle
  (0.4141,0.5548) -- (0.4297,0.5548) -- (0.4219,0.5683) -- cycle (0.3906,0.5683) -- (0.4062,0.5683) -- (0.3984,0.5819) -- cycle (0.4062,0.5683) -- (0.4219,0.5683) -- (0.4141,0.5819) -- cycle (0.3984,0.5819) -- (0.4141,0.5819) -- (0.4062,0.5954) -- cycle
  (0.3438,0.5954) -- (0.3594,0.5954) -- (0.3516,0.6089) -- cycle (0.3594,0.5954) -- (0.375,0.5954) -- (0.3672,0.6089) -- cycle (0.3516,0.6089) -- (0.3672,0.6089) -- (0.3594,0.6225) -- cycle (0.375,0.5954) -- (0.3906,0.5954) -- (0.3828,0.6089) -- cycle
  (0.3906,0.5954) -- (0.4062,0.5954) -- (0.3984,0.6089) -- cycle (0.3828,0.6089) -- (0.3984,0.6089) -- (0.3906,0.6225) -- cycle (0.3594,0.6225) -- (0.375,0.6225) -- (0.3672,0.636) -- cycle (0.375,0.6225) -- (0.3906,0.6225) -- (0.3828,0.636) -- cycle
  (0.3672,0.636) -- (0.3828,0.636) -- (0.375,0.6495) -- cycle (0.5,0.433) -- (0.5156,0.433) -- (0.5078,0.4465) -- cycle (0.5156,0.433) -- (0.5312,0.433) -- (0.5234,0.4465) -- cycle (0.5078,0.4465) -- (0.5234,0.4465) -- (0.5156,0.4601) -- cycle
  (0.5312,0.433) -- (0.5469,0.433) -- (0.5391,0.4465) -- cycle (0.5469,0.433) -- (0.5625,0.433) -- (0.5547,0.4465) -- cycle (0.5391,0.4465) -- (0.5547,0.4465) -- (0.5469,0.4601) -- cycle (0.5156,0.4601) -- (0.5312,0.4601) -- (0.5234,0.4736) -- cycle
  (0.5312,0.4601) -- (0.5469,0.4601) -- (0.5391,0.4736) -- cycle (0.5234,0.4736) -- (0.5391,0.4736) -- (0.5312,0.4871) -- cycle (0.5625,0.433) -- (0.5781,0.433) -- (0.5703,0.4465) -- cycle (0.5781,0.433) -- (0.5938,0.433) -- (0.5859,0.4465) -- cycle
  (0.5703,0.4465) -- (0.5859,0.4465) -- (0.5781,0.4601) -- cycle (0.5938,0.433) -- (0.6094,0.433) -- (0.6016,0.4465) -- cycle (0.6094,0.433) -- (0.625,0.433) -- (0.6172,0.4465) -- cycle (0.6016,0.4465) -- (0.6172,0.4465) -- (0.6094,0.4601) -- cycle
  (0.5781,0.4601) -- (0.5938,0.4601) -- (0.5859,0.4736) -- cycle (0.5938,0.4601) -- (0.6094,0.4601) -- (0.6016,0.4736) -- cycle (0.5859,0.4736) -- (0.6016,0.4736) -- (0.5938,0.4871) -- cycle (0.5312,0.4871) -- (0.5469,0.4871) -- (0.5391,0.5007) -- cycle
  (0.5469,0.4871) -- (0.5625,0.4871) -- (0.5547,0.5007) -- cycle (0.5391,0.5007) -- (0.5547,0.5007) -- (0.5469,0.5142) -- cycle (0.5625,0.4871) -- (0.5781,0.4871) -- (0.5703,0.5007) -- cycle (0.5781,0.4871) -- (0.5938,0.4871) -- (0.5859,0.5007) -- cycle
  (0.5703,0.5007) -- (0.5859,0.5007) -- (0.5781,0.5142) -- cycle (0.5469,0.5142) -- (0.5625,0.5142) -- (0.5547,0.5277) -- cycle (0.5625,0.5142) -- (0.5781,0.5142) -- (0.5703,0.5277) -- cycle (0.5547,0.5277) -- (0.5703,0.5277) -- (0.5625,0.5413) -- cycle
  (0.625,0.433) -- (0.6406,0.433) -- (0.6328,0.4465) -- cycle (0.6406,0.433) -- (0.6562,0.433) -- (0.6484,0.4465) -- cycle (0.6328,0.4465) -- (0.6484,0.4465) -- (0.6406,0.4601) -- cycle (0.6562,0.433) -- (0.6719,0.433) -- (0.6641,0.4465) -- cycle
  (0.6719,0.433) -- (0.6875,0.433) -- (0.6797,0.4465) -- cycle (0.6641,0.4465) -- (0.6797,0.4465) -- (0.6719,0.4601) -- cycle (0.6406,0.4601) -- (0.6562,0.4601) -- (0.6484,0.4736) -- cycle (0.6562,0.4601) -- (0.6719,0.4601) -- (0.6641,0.4736) -- cycle
  (0.6484,0.4736) -- (0.6641,0.4736) -- (0.6562,0.4871) -- cycle (0.6875,0.433) -- (0.7031,0.433) -- (0.6953,0.4465) -- cycle (0.7031,0.433) -- (0.7188,0.433) -- (0.7109,0.4465) -- cycle (0.6953,0.4465) -- (0.7109,0.4465) -- (0.7031,0.4601) -- cycle
  (0.7188,0.433) -- (0.7344,0.433) -- (0.7266,0.4465) -- cycle (0.7344,0.433) -- (0.75,0.433) -- (0.7422,0.4465) -- cycle (0.7266,0.4465) -- (0.7422,0.4465) -- (0.7344,0.4601) -- cycle (0.7031,0.4601) -- (0.7188,0.4601) -- (0.7109,0.4736) -- cycle
  (0.7188,0.4601) -- (0.7344,0.4601) -- (0.7266,0.4736) -- cycle (0.7109,0.4736) -- (0.7266,0.4736) -- (0.7188,0.4871) -- cycle (0.6562,0.4871) -- (0.6719,0.4871) -- (0.6641,0.5007) -- cycle (0.6719,0.4871) -- (0.6875,0.4871) -- (0.6797,0.5007) -- cycle
  (0.6641,0.5007) -- (0.6797,0.5007) -- (0.6719,0.5142) -- cycle (0.6875,0.4871) -- (0.7031,0.4871) -- (0.6953,0.5007) -- cycle (0.7031,0.4871) -- (0.7188,0.4871) -- (0.7109,0.5007) -- cycle (0.6953,0.5007) -- (0.7109,0.5007) -- (0.7031,0.5142) -- cycle
  (0.6719,0.5142) -- (0.6875,0.5142) -- (0.6797,0.5277) -- cycle (0.6875,0.5142) -- (0.7031,0.5142) -- (0.6953,0.5277) -- cycle (0.6797,0.5277) -- (0.6953,0.5277) -- (0.6875,0.5413) -- cycle (0.5625,0.5413) -- (0.5781,0.5413) -- (0.5703,0.5548) -- cycle
  (0.5781,0.5413) -- (0.5938,0.5413) -- (0.5859,0.5548) -- cycle (0.5703,0.5548) -- (0.5859,0.5548) -- (0.5781,0.5683) -- cycle (0.5938,0.5413) -- (0.6094,0.5413) -- (0.6016,0.5548) -- cycle (0.6094,0.5413) -- (0.625,0.5413) -- (0.6172,0.5548) -- cycle
  (0.6016,0.5548) -- (0.6172,0.5548) -- (0.6094,0.5683) -- cycle (0.5781,0.5683) -- (0.5938,0.5683) -- (0.5859,0.5819) -- cycle (0.5938,0.5683) -- (0.6094,0.5683) -- (0.6016,0.5819) -- cycle (0.5859,0.5819) -- (0.6016,0.5819) -- (0.5938,0.5954) -- cycle
  (0.625,0.5413) -- (0.6406,0.5413) -- (0.6328,0.5548) -- cycle (0.6406,0.5413) -- (0.6562,0.5413) -- (0.6484,0.5548) -- cycle (0.6328,0.5548) -- (0.6484,0.5548) -- (0.6406,0.5683) -- cycle (0.6562,0.5413) -- (0.6719,0.5413) -- (0.6641,0.5548) -- cycle
  (0.6719,0.5413) -- (0.6875,0.5413) -- (0.6797,0.5548) -- cycle (0.6641,0.5548) -- (0.6797,0.5548) -- (0.6719,0.5683) -- cycle (0.6406,0.5683) -- (0.6562,0.5683) -- (0.6484,0.5819) -- cycle (0.6562,0.5683) -- (0.6719,0.5683) -- (0.6641,0.5819) -- cycle
  (0.6484,0.5819) -- (0.6641,0.5819) -- (0.6562,0.5954) -- cycle (0.5938,0.5954) -- (0.6094,0.5954) -- (0.6016,0.6089) -- cycle (0.6094,0.5954) -- (0.625,0.5954) -- (0.6172,0.6089) -- cycle (0.6016,0.6089) -- (0.6172,0.6089) -- (0.6094,0.6225) -- cycle
  (0.625,0.5954) -- (0.6406,0.5954) -- (0.6328,0.6089) -- cycle (0.6406,0.5954) -- (0.6562,0.5954) -- (0.6484,0.6089) -- cycle (0.6328,0.6089) -- (0.6484,0.6089) -- (0.6406,0.6225) -- cycle (0.6094,0.6225) -- (0.625,0.6225) -- (0.6172,0.636) -- cycle
  (0.625,0.6225) -- (0.6406,0.6225) -- (0.6328,0.636) -- cycle (0.6172,0.636) -- (0.6328,0.636) -- (0.625,0.6495) -- cycle (0.375,0.6495) -- (0.3906,0.6495) -- (0.3828,0.6631) -- cycle (0.3906,0.6495) -- (0.4062,0.6495) -- (0.3984,0.6631) -- cycle
  (0.3828,0.6631) -- (0.3984,0.6631) -- (0.3906,0.6766) -- cycle (0.4062,0.6495) -- (0.4219,0.6495) -- (0.4141,0.6631) -- cycle (0.4219,0.6495) -- (0.4375,0.6495) -- (0.4297,0.6631) -- cycle (0.4141,0.6631) -- (0.4297,0.6631) -- (0.4219,0.6766) -- cycle
  (0.3906,0.6766) -- (0.4062,0.6766) -- (0.3984,0.6901) -- cycle (0.4062,0.6766) -- (0.4219,0.6766) -- (0.4141,0.6901) -- cycle (0.3984,0.6901) -- (0.4141,0.6901) -- (0.4062,0.7036) -- cycle (0.4375,0.6495) -- (0.4531,0.6495) -- (0.4453,0.6631) -- cycle
  (0.4531,0.6495) -- (0.4688,0.6495) -- (0.4609,0.6631) -- cycle (0.4453,0.6631) -- (0.4609,0.6631) -- (0.4531,0.6766) -- cycle (0.4688,0.6495) -- (0.4844,0.6495) -- (0.4766,0.6631) -- cycle (0.4844,0.6495) -- (0.5,0.6495) -- (0.4922,0.6631) -- cycle
  (0.4766,0.6631) -- (0.4922,0.6631) -- (0.4844,0.6766) -- cycle (0.4531,0.6766) -- (0.4688,0.6766) -- (0.4609,0.6901) -- cycle (0.4688,0.6766) -- (0.4844,0.6766) -- (0.4766,0.6901) -- cycle (0.4609,0.6901) -- (0.4766,0.6901) -- (0.4688,0.7036) -- cycle
  (0.4062,0.7036) -- (0.4219,0.7036) -- (0.4141,0.7172) -- cycle (0.4219,0.7036) -- (0.4375,0.7036) -- (0.4297,0.7172) -- cycle (0.4141,0.7172) -- (0.4297,0.7172) -- (0.4219,0.7307) -- cycle (0.4375,0.7036) -- (0.4531,0.7036) -- (0.4453,0.7172) -- cycle
  (0.4531,0.7036) -- (0.4688,0.7036) -- (0.4609,0.7172) -- cycle (0.4453,0.7172) -- (0.4609,0.7172) -- (0.4531,0.7307) -- cycle (0.4219,0.7307) -- (0.4375,0.7307) -- (0.4297,0.7442) -- cycle (0.4375,0.7307) -- (0.4531,0.7307) -- (0.4453,0.7442) -- cycle
  (0.4297,0.7442) -- (0.4453,0.7442) -- (0.4375,0.7578) -- cycle (0.5,0.6495) -- (0.5156,0.6495) -- (0.5078,0.6631) -- cycle (0.5156,0.6495) -- (0.5312,0.6495) -- (0.5234,0.6631) -- cycle (0.5078,0.6631) -- (0.5234,0.6631) -- (0.5156,0.6766) -- cycle
  (0.5312,0.6495) -- (0.5469,0.6495) -- (0.5391,0.6631) -- cycle (0.5469,0.6495) -- (0.5625,0.6495) -- (0.5547,0.6631) -- cycle (0.5391,0.6631) -- (0.5547,0.6631) -- (0.5469,0.6766) -- cycle (0.5156,0.6766) -- (0.5312,0.6766) -- (0.5234,0.6901) -- cycle
  (0.5312,0.6766) -- (0.5469,0.6766) -- (0.5391,0.6901) -- cycle (0.5234,0.6901) -- (0.5391,0.6901) -- (0.5312,0.7036) -- cycle (0.5625,0.6495) -- (0.5781,0.6495) -- (0.5703,0.6631) -- cycle (0.5781,0.6495) -- (0.5938,0.6495) -- (0.5859,0.6631) -- cycle
  (0.5703,0.6631) -- (0.5859,0.6631) -- (0.5781,0.6766) -- cycle (0.5938,0.6495) -- (0.6094,0.6495) -- (0.6016,0.6631) -- cycle (0.6094,0.6495) -- (0.625,0.6495) -- (0.6172,0.6631) -- cycle (0.6016,0.6631) -- (0.6172,0.6631) -- (0.6094,0.6766) -- cycle
  (0.5781,0.6766) -- (0.5938,0.6766) -- (0.5859,0.6901) -- cycle (0.5938,0.6766) -- (0.6094,0.6766) -- (0.6016,0.6901) -- cycle (0.5859,0.6901) -- (0.6016,0.6901) -- (0.5938,0.7036) -- cycle (0.5312,0.7036) -- (0.5469,0.7036) -- (0.5391,0.7172) -- cycle
  (0.5469,0.7036) -- (0.5625,0.7036) -- (0.5547,0.7172) -- cycle (0.5391,0.7172) -- (0.5547,0.7172) -- (0.5469,0.7307) -- cycle (0.5625,0.7036) -- (0.5781,0.7036) -- (0.5703,0.7172) -- cycle (0.5781,0.7036) -- (0.5938,0.7036) -- (0.5859,0.7172) -- cycle
  (0.5703,0.7172) -- (0.5859,0.7172) -- (0.5781,0.7307) -- cycle (0.5469,0.7307) -- (0.5625,0.7307) -- (0.5547,0.7442) -- cycle (0.5625,0.7307) -- (0.5781,0.7307) -- (0.5703,0.7442) -- cycle (0.5547,0.7442) -- (0.5703,0.7442) -- (0.5625,0.7578) -- cycle
  (0.4375,0.7578) -- (0.4531,0.7578) -- (0.4453,0.7713) -- cycle (0.4531,0.7578) -- (0.4688,0.7578) -- (0.4609,0.7713) -- cycle (0.4453,0.7713) -- (0.4609,0.7713) -- (0.4531,0.7848) -- cycle (0.4688,0.7578) -- (0.4844,0.7578) -- (0.4766,0.7713) -- cycle
  (0.4844,0.7578) -- (0.5,0.7578) -- (0.4922,0.7713) -- cycle (0.4766,0.7713) -- (0.4922,0.7713) -- (0.4844,0.7848) -- cycle (0.4531,0.7848) -- (0.4688,0.7848) -- (0.4609,0.7984) -- cycle (0.4688,0.7848) -- (0.4844,0.7848) -- (0.4766,0.7984) -- cycle
  (0.4609,0.7984) -- (0.4766,0.7984) -- (0.4688,0.8119) -- cycle (0.5,0.7578) -- (0.5156,0.7578) -- (0.5078,0.7713) -- cycle (0.5156,0.7578) -- (0.5312,0.7578) -- (0.5234,0.7713) -- cycle (0.5078,0.7713) -- (0.5234,0.7713) -- (0.5156,0.7848) -- cycle
  (0.5312,0.7578) -- (0.5469,0.7578) -- (0.5391,0.7713) -- cycle (0.5469,0.7578) -- (0.5625,0.7578) -- (0.5547,0.7713) -- cycle (0.5391,0.7713) -- (0.5547,0.7713) -- (0.5469,0.7848) -- cycle (0.5156,0.7848) -- (0.5312,0.7848) -- (0.5234,0.7984) -- cycle
  (0.5312,0.7848) -- (0.5469,0.7848) -- (0.5391,0.7984) -- cycle (0.5234,0.7984) -- (0.5391,0.7984) -- (0.5312,0.8119) -- cycle (0.4688,0.8119) -- (0.4844,0.8119) -- (0.4766,0.8254) -- cycle (0.4844,0.8119) -- (0.5,0.8119) -- (0.4922,0.8254) -- cycle
  (0.4766,0.8254) -- (0.4922,0.8254) -- (0.4844,0.839) -- cycle (0.5,0.8119) -- (0.5156,0.8119) -- (0.5078,0.8254) -- cycle (0.5156,0.8119) -- (0.5312,0.8119) -- (0.5234,0.8254) -- cycle (0.5078,0.8254) -- (0.5234,0.8254) -- (0.5156,0.839) -- cycle
  (0.4844,0.839) -- (0.5,0.839) -- (0.4922,0.8525) -- cycle (0.5,0.839) -- (0.5156,0.839) -- (0.5078,0.8525) -- cycle (0.4922,0.8525) -- (0.5078,0.8525) -- (0.5,0.866) -- cycle;

      \node[etiq, firebrick3, below] at (0.25,-0.02) {$w_1(S)$};
      \node[etiq, steelblue4, below] at (0.75,-0.02) {$w_2(S)$};
      \node[etiq, coveraccent!80!black, right] at (0.68,0.62) {$w_3(S)$};
      \begin{scope}[shift={(1.3,0.18)}, scale=1.35]
        % Courbe de von Koch : les quatre copies w_i(K)
% Fichier généré par generer.py : ne pas modifier à la main.
\draw[trait1]
  (0,0) -- (0.0041,0) -- (0.0062,0.0036) -- (0.0082,0) -- (0.0123,0) -- (0.0144,0.0036) --
  (0.0123,0.0071) -- (0.0165,0.0071) -- (0.0185,0.0107) -- (0.0206,0.0071) -- (0.0247,0.0071) -- (0.0226,0.0036) --
  (0.0247,0) -- (0.0288,0) -- (0.0309,0.0036) -- (0.0329,0) -- (0.037,0) -- (0.0391,0.0036) --
  (0.037,0.0071) -- (0.0412,0.0071) -- (0.0432,0.0107) -- (0.0412,0.0143) -- (0.037,0.0143) -- (0.0391,0.0178) --
  (0.037,0.0214) -- (0.0412,0.0214) -- (0.0432,0.0249) -- (0.0453,0.0214) -- (0.0494,0.0214) -- (0.0514,0.0249) --
  (0.0494,0.0285) -- (0.0535,0.0285) -- (0.0556,0.0321) -- (0.0576,0.0285) -- (0.0617,0.0285) -- (0.0597,0.0249) --
  (0.0617,0.0214) -- (0.0658,0.0214) -- (0.0679,0.0249) -- (0.07,0.0214) -- (0.0741,0.0214) -- (0.072,0.0178) --
  (0.0741,0.0143) -- (0.07,0.0143) -- (0.0679,0.0107) -- (0.07,0.0071) -- (0.0741,0.0071) -- (0.072,0.0036) --
  (0.0741,0) -- (0.0782,0) -- (0.0802,0.0036) -- (0.0823,0) -- (0.0864,0) -- (0.0885,0.0036) --
  (0.0864,0.0071) -- (0.0905,0.0071) -- (0.0926,0.0107) -- (0.0947,0.0071) -- (0.0988,0.0071) -- (0.0967,0.0036) --
  (0.0988,0) -- (0.1029,0) -- (0.1049,0.0036) -- (0.107,0) -- (0.1111,0) -- (0.1132,0.0036) --
  (0.1111,0.0071) -- (0.1152,0.0071) -- (0.1173,0.0107) -- (0.1152,0.0143) -- (0.1111,0.0143) -- (0.1132,0.0178) --
  (0.1111,0.0214) -- (0.1152,0.0214) -- (0.1173,0.0249) -- (0.1193,0.0214) -- (0.1235,0.0214) -- (0.1255,0.0249) --
  (0.1235,0.0285) -- (0.1276,0.0285) -- (0.1296,0.0321) -- (0.1276,0.0356) -- (0.1235,0.0356) -- (0.1255,0.0392) --
  (0.1235,0.0428) -- (0.1193,0.0428) -- (0.1173,0.0392) -- (0.1152,0.0428) -- (0.1111,0.0428) -- (0.1132,0.0463) --
  (0.1111,0.0499) -- (0.1152,0.0499) -- (0.1173,0.0535) -- (0.1152,0.057) -- (0.1111,0.057) -- (0.1132,0.0606) --
  (0.1111,0.0642) -- (0.1152,0.0642) -- (0.1173,0.0677) -- (0.1193,0.0642) -- (0.1235,0.0642) -- (0.1255,0.0677) --
  (0.1235,0.0713) -- (0.1276,0.0713) -- (0.1296,0.0748) -- (0.1317,0.0713) -- (0.1358,0.0713) -- (0.1337,0.0677) --
  (0.1358,0.0642) -- (0.1399,0.0642) -- (0.142,0.0677) -- (0.144,0.0642) -- (0.1481,0.0642) -- (0.1502,0.0677) --
  (0.1481,0.0713) -- (0.1523,0.0713) -- (0.1543,0.0748) -- (0.1523,0.0784) -- (0.1481,0.0784) -- (0.1502,0.082) --
  (0.1481,0.0855) -- (0.1523,0.0855) -- (0.1543,0.0891) -- (0.1564,0.0855) -- (0.1605,0.0855) -- (0.1626,0.0891) --
  (0.1605,0.0927) -- (0.1646,0.0927) -- (0.1667,0.0962) -- (0.1687,0.0927) -- (0.1728,0.0927) -- (0.1708,0.0891) --
  (0.1728,0.0855) -- (0.177,0.0855) -- (0.179,0.0891) -- (0.1811,0.0855) -- (0.1852,0.0855) -- (0.1831,0.082) --
  (0.1852,0.0784) -- (0.1811,0.0784) -- (0.179,0.0748) -- (0.1811,0.0713) -- (0.1852,0.0713) -- (0.1831,0.0677) --
  (0.1852,0.0642) -- (0.1893,0.0642) -- (0.1914,0.0677) -- (0.1934,0.0642) -- (0.1975,0.0642) -- (0.1996,0.0677) --
  (0.1975,0.0713) -- (0.2016,0.0713) -- (0.2037,0.0748) -- (0.2058,0.0713) -- (0.2099,0.0713) -- (0.2078,0.0677) --
  (0.2099,0.0642) -- (0.214,0.0642) -- (0.216,0.0677) -- (0.2181,0.0642) -- (0.2222,0.0642) -- (0.2202,0.0606) --
  (0.2222,0.057) -- (0.2181,0.057) -- (0.216,0.0535) -- (0.2181,0.0499) -- (0.2222,0.0499) -- (0.2202,0.0463) --
  (0.2222,0.0428) -- (0.2181,0.0428) -- (0.216,0.0392) -- (0.214,0.0428) -- (0.2099,0.0428) -- (0.2078,0.0392) --
  (0.2099,0.0356) -- (0.2058,0.0356) -- (0.2037,0.0321) -- (0.2058,0.0285) -- (0.2099,0.0285) -- (0.2078,0.0249) --
  (0.2099,0.0214) -- (0.214,0.0214) -- (0.216,0.0249) -- (0.2181,0.0214) -- (0.2222,0.0214) -- (0.2202,0.0178) --
  (0.2222,0.0143) -- (0.2181,0.0143) -- (0.216,0.0107) -- (0.2181,0.0071) -- (0.2222,0.0071) -- (0.2202,0.0036) --
  (0.2222,0) -- (0.2263,0) -- (0.2284,0.0036) -- (0.2305,0) -- (0.2346,0) -- (0.2366,0.0036) --
  (0.2346,0.0071) -- (0.2387,0.0071) -- (0.2407,0.0107) -- (0.2428,0.0071) -- (0.2469,0.0071) -- (0.2449,0.0036) --
  (0.2469,0) -- (0.251,0) -- (0.2531,0.0036) -- (0.2551,0) -- (0.2593,0) -- (0.2613,0.0036) --
  (0.2593,0.0071) -- (0.2634,0.0071) -- (0.2654,0.0107) -- (0.2634,0.0143) -- (0.2593,0.0143) -- (0.2613,0.0178) --
  (0.2593,0.0214) -- (0.2634,0.0214) -- (0.2654,0.0249) -- (0.2675,0.0214) -- (0.2716,0.0214) -- (0.2737,0.0249) --
  (0.2716,0.0285) -- (0.2757,0.0285) -- (0.2778,0.0321) -- (0.2798,0.0285) -- (0.284,0.0285) -- (0.2819,0.0249) --
  (0.284,0.0214) -- (0.2881,0.0214) -- (0.2901,0.0249) -- (0.2922,0.0214) -- (0.2963,0.0214) -- (0.2942,0.0178) --
  (0.2963,0.0143) -- (0.2922,0.0143) -- (0.2901,0.0107) -- (0.2922,0.0071) -- (0.2963,0.0071) -- (0.2942,0.0036) --
  (0.2963,0) -- (0.3004,0) -- (0.3025,0.0036) -- (0.3045,0) -- (0.3086,0) -- (0.3107,0.0036) --
  (0.3086,0.0071) -- (0.3128,0.0071) -- (0.3148,0.0107) -- (0.3169,0.0071) -- (0.321,0.0071) -- (0.3189,0.0036) --
  (0.321,0) -- (0.3251,0) -- (0.3272,0.0036) -- (0.3292,0) -- (0.3333,0);
\draw[trait2]
  (0.3333,0) -- (0.3354,0.0036) -- (0.3333,0.0071) -- (0.3374,0.0071) -- (0.3395,0.0107) -- (0.3374,0.0143) --
  (0.3333,0.0143) -- (0.3354,0.0178) -- (0.3333,0.0214) -- (0.3374,0.0214) -- (0.3395,0.0249) -- (0.3416,0.0214) --
  (0.3457,0.0214) -- (0.3477,0.0249) -- (0.3457,0.0285) -- (0.3498,0.0285) -- (0.3519,0.0321) -- (0.3498,0.0356) --
  (0.3457,0.0356) -- (0.3477,0.0392) -- (0.3457,0.0428) -- (0.3416,0.0428) -- (0.3395,0.0392) -- (0.3374,0.0428) --
  (0.3333,0.0428) -- (0.3354,0.0463) -- (0.3333,0.0499) -- (0.3374,0.0499) -- (0.3395,0.0535) -- (0.3374,0.057) --
  (0.3333,0.057) -- (0.3354,0.0606) -- (0.3333,0.0642) -- (0.3374,0.0642) -- (0.3395,0.0677) -- (0.3416,0.0642) --
  (0.3457,0.0642) -- (0.3477,0.0677) -- (0.3457,0.0713) -- (0.3498,0.0713) -- (0.3519,0.0748) -- (0.3539,0.0713) --
  (0.358,0.0713) -- (0.356,0.0677) -- (0.358,0.0642) -- (0.3621,0.0642) -- (0.3642,0.0677) -- (0.3663,0.0642) --
  (0.3704,0.0642) -- (0.3724,0.0677) -- (0.3704,0.0713) -- (0.3745,0.0713) -- (0.3765,0.0748) -- (0.3745,0.0784) --
  (0.3704,0.0784) -- (0.3724,0.082) -- (0.3704,0.0855) -- (0.3745,0.0855) -- (0.3765,0.0891) -- (0.3786,0.0855) --
  (0.3827,0.0855) -- (0.3848,0.0891) -- (0.3827,0.0927) -- (0.3868,0.0927) -- (0.3889,0.0962) -- (0.3868,0.0998) --
  (0.3827,0.0998) -- (0.3848,0.1034) -- (0.3827,0.1069) -- (0.3786,0.1069) -- (0.3765,0.1034) -- (0.3745,0.1069) --
  (0.3704,0.1069) -- (0.3724,0.1105) -- (0.3704,0.114) -- (0.3745,0.114) -- (0.3765,0.1176) -- (0.3745,0.1212) --
  (0.3704,0.1212) -- (0.3724,0.1247) -- (0.3704,0.1283) -- (0.3663,0.1283) -- (0.3642,0.1247) -- (0.3621,0.1283) --
  (0.358,0.1283) -- (0.356,0.1247) -- (0.358,0.1212) -- (0.3539,0.1212) -- (0.3519,0.1176) -- (0.3498,0.1212) --
  (0.3457,0.1212) -- (0.3477,0.1247) -- (0.3457,0.1283) -- (0.3416,0.1283) -- (0.3395,0.1247) -- (0.3374,0.1283) --
  (0.3333,0.1283) -- (0.3354,0.1319) -- (0.3333,0.1354) -- (0.3374,0.1354) -- (0.3395,0.139) -- (0.3374,0.1426) --
  (0.3333,0.1426) -- (0.3354,0.1461) -- (0.3333,0.1497) -- (0.3374,0.1497) -- (0.3395,0.1532) -- (0.3416,0.1497) --
  (0.3457,0.1497) -- (0.3477,0.1532) -- (0.3457,0.1568) -- (0.3498,0.1568) -- (0.3519,0.1604) -- (0.3498,0.1639) --
  (0.3457,0.1639) -- (0.3477,0.1675) -- (0.3457,0.1711) -- (0.3416,0.1711) -- (0.3395,0.1675) -- (0.3374,0.1711) --
  (0.3333,0.1711) -- (0.3354,0.1746) -- (0.3333,0.1782) -- (0.3374,0.1782) -- (0.3395,0.1818) -- (0.3374,0.1853) --
  (0.3333,0.1853) -- (0.3354,0.1889) -- (0.3333,0.1925) -- (0.3374,0.1925) -- (0.3395,0.196) -- (0.3416,0.1925) --
  (0.3457,0.1925) -- (0.3477,0.196) -- (0.3457,0.1996) -- (0.3498,0.1996) -- (0.3519,0.2031) -- (0.3539,0.1996) --
  (0.358,0.1996) -- (0.356,0.196) -- (0.358,0.1925) -- (0.3621,0.1925) -- (0.3642,0.196) -- (0.3663,0.1925) --
  (0.3704,0.1925) -- (0.3724,0.196) -- (0.3704,0.1996) -- (0.3745,0.1996) -- (0.3765,0.2031) -- (0.3745,0.2067) --
  (0.3704,0.2067) -- (0.3724,0.2103) -- (0.3704,0.2138) -- (0.3745,0.2138) -- (0.3765,0.2174) -- (0.3786,0.2138) --
  (0.3827,0.2138) -- (0.3848,0.2174) -- (0.3827,0.221) -- (0.3868,0.221) -- (0.3889,0.2245) -- (0.3909,0.221) --
  (0.3951,0.221) -- (0.393,0.2174) -- (0.3951,0.2138) -- (0.3992,0.2138) -- (0.4012,0.2174) -- (0.4033,0.2138) --
  (0.4074,0.2138) -- (0.4053,0.2103) -- (0.4074,0.2067) -- (0.4033,0.2067) -- (0.4012,0.2031) -- (0.4033,0.1996) --
  (0.4074,0.1996) -- (0.4053,0.196) -- (0.4074,0.1925) -- (0.4115,0.1925) -- (0.4136,0.196) -- (0.4156,0.1925) --
  (0.4198,0.1925) -- (0.4218,0.196) -- (0.4198,0.1996) -- (0.4239,0.1996) -- (0.4259,0.2031) -- (0.428,0.1996) --
  (0.4321,0.1996) -- (0.43,0.196) -- (0.4321,0.1925) -- (0.4362,0.1925) -- (0.4383,0.196) -- (0.4403,0.1925) --
  (0.4444,0.1925) -- (0.4465,0.196) -- (0.4444,0.1996) -- (0.4486,0.1996) -- (0.4506,0.2031) -- (0.4486,0.2067) --
  (0.4444,0.2067) -- (0.4465,0.2103) -- (0.4444,0.2138) -- (0.4486,0.2138) -- (0.4506,0.2174) -- (0.4527,0.2138) --
  (0.4568,0.2138) -- (0.4588,0.2174) -- (0.4568,0.221) -- (0.4609,0.221) -- (0.463,0.2245) -- (0.4609,0.2281) --
  (0.4568,0.2281) -- (0.4588,0.2317) -- (0.4568,0.2352) -- (0.4527,0.2352) -- (0.4506,0.2317) -- (0.4486,0.2352) --
  (0.4444,0.2352) -- (0.4465,0.2388) -- (0.4444,0.2423) -- (0.4486,0.2423) -- (0.4506,0.2459) -- (0.4486,0.2495) --
  (0.4444,0.2495) -- (0.4465,0.253) -- (0.4444,0.2566) -- (0.4486,0.2566) -- (0.4506,0.2602) -- (0.4527,0.2566) --
  (0.4568,0.2566) -- (0.4588,0.2602) -- (0.4568,0.2637) -- (0.4609,0.2637) -- (0.463,0.2673) -- (0.465,0.2637) --
  (0.4691,0.2637) -- (0.4671,0.2602) -- (0.4691,0.2566) -- (0.4733,0.2566) -- (0.4753,0.2602) -- (0.4774,0.2566) --
  (0.4815,0.2566) -- (0.4835,0.2602) -- (0.4815,0.2637) -- (0.4856,0.2637) -- (0.4877,0.2673) -- (0.4856,0.2709) --
  (0.4815,0.2709) -- (0.4835,0.2744) -- (0.4815,0.278) -- (0.4856,0.278) -- (0.4877,0.2815) -- (0.4897,0.278) --
  (0.4938,0.278) -- (0.4959,0.2815) -- (0.4938,0.2851) -- (0.4979,0.2851) -- (0.5,0.2887);
\draw[trait3]
  (0.5,0.2887) -- (0.5021,0.2851) -- (0.5062,0.2851) -- (0.5041,0.2815) -- (0.5062,0.278) -- (0.5103,0.278) --
  (0.5123,0.2815) -- (0.5144,0.278) -- (0.5185,0.278) -- (0.5165,0.2744) -- (0.5185,0.2709) -- (0.5144,0.2709) --
  (0.5123,0.2673) -- (0.5144,0.2637) -- (0.5185,0.2637) -- (0.5165,0.2602) -- (0.5185,0.2566) -- (0.5226,0.2566) --
  (0.5247,0.2602) -- (0.5267,0.2566) -- (0.5309,0.2566) -- (0.5329,0.2602) -- (0.5309,0.2637) -- (0.535,0.2637) --
  (0.537,0.2673) -- (0.5391,0.2637) -- (0.5432,0.2637) -- (0.5412,0.2602) -- (0.5432,0.2566) -- (0.5473,0.2566) --
  (0.5494,0.2602) -- (0.5514,0.2566) -- (0.5556,0.2566) -- (0.5535,0.253) -- (0.5556,0.2495) -- (0.5514,0.2495) --
  (0.5494,0.2459) -- (0.5514,0.2423) -- (0.5556,0.2423) -- (0.5535,0.2388) -- (0.5556,0.2352) -- (0.5514,0.2352) --
  (0.5494,0.2317) -- (0.5473,0.2352) -- (0.5432,0.2352) -- (0.5412,0.2317) -- (0.5432,0.2281) -- (0.5391,0.2281) --
  (0.537,0.2245) -- (0.5391,0.221) -- (0.5432,0.221) -- (0.5412,0.2174) -- (0.5432,0.2138) -- (0.5473,0.2138) --
  (0.5494,0.2174) -- (0.5514,0.2138) -- (0.5556,0.2138) -- (0.5535,0.2103) -- (0.5556,0.2067) -- (0.5514,0.2067) --
  (0.5494,0.2031) -- (0.5514,0.1996) -- (0.5556,0.1996) -- (0.5535,0.196) -- (0.5556,0.1925) -- (0.5597,0.1925) --
  (0.5617,0.196) -- (0.5638,0.1925) -- (0.5679,0.1925) -- (0.57,0.196) -- (0.5679,0.1996) -- (0.572,0.1996) --
  (0.5741,0.2031) -- (0.5761,0.1996) -- (0.5802,0.1996) -- (0.5782,0.196) -- (0.5802,0.1925) -- (0.5844,0.1925) --
  (0.5864,0.196) -- (0.5885,0.1925) -- (0.5926,0.1925) -- (0.5947,0.196) -- (0.5926,0.1996) -- (0.5967,0.1996) --
  (0.5988,0.2031) -- (0.5967,0.2067) -- (0.5926,0.2067) -- (0.5947,0.2103) -- (0.5926,0.2138) -- (0.5967,0.2138) --
  (0.5988,0.2174) -- (0.6008,0.2138) -- (0.6049,0.2138) -- (0.607,0.2174) -- (0.6049,0.221) -- (0.6091,0.221) --
  (0.6111,0.2245) -- (0.6132,0.221) -- (0.6173,0.221) -- (0.6152,0.2174) -- (0.6173,0.2138) -- (0.6214,0.2138) --
  (0.6235,0.2174) -- (0.6255,0.2138) -- (0.6296,0.2138) -- (0.6276,0.2103) -- (0.6296,0.2067) -- (0.6255,0.2067) --
  (0.6235,0.2031) -- (0.6255,0.1996) -- (0.6296,0.1996) -- (0.6276,0.196) -- (0.6296,0.1925) -- (0.6337,0.1925) --
  (0.6358,0.196) -- (0.6379,0.1925) -- (0.642,0.1925) -- (0.644,0.196) -- (0.642,0.1996) -- (0.6461,0.1996) --
  (0.6481,0.2031) -- (0.6502,0.1996) -- (0.6543,0.1996) -- (0.6523,0.196) -- (0.6543,0.1925) -- (0.6584,0.1925) --
  (0.6605,0.196) -- (0.6626,0.1925) -- (0.6667,0.1925) -- (0.6646,0.1889) -- (0.6667,0.1853) -- (0.6626,0.1853) --
  (0.6605,0.1818) -- (0.6626,0.1782) -- (0.6667,0.1782) -- (0.6646,0.1746) -- (0.6667,0.1711) -- (0.6626,0.1711) --
  (0.6605,0.1675) -- (0.6584,0.1711) -- (0.6543,0.1711) -- (0.6523,0.1675) -- (0.6543,0.1639) -- (0.6502,0.1639) --
  (0.6481,0.1604) -- (0.6502,0.1568) -- (0.6543,0.1568) -- (0.6523,0.1532) -- (0.6543,0.1497) -- (0.6584,0.1497) --
  (0.6605,0.1532) -- (0.6626,0.1497) -- (0.6667,0.1497) -- (0.6646,0.1461) -- (0.6667,0.1426) -- (0.6626,0.1426) --
  (0.6605,0.139) -- (0.6626,0.1354) -- (0.6667,0.1354) -- (0.6646,0.1319) -- (0.6667,0.1283) -- (0.6626,0.1283) --
  (0.6605,0.1247) -- (0.6584,0.1283) -- (0.6543,0.1283) -- (0.6523,0.1247) -- (0.6543,0.1212) -- (0.6502,0.1212) --
  (0.6481,0.1176) -- (0.6461,0.1212) -- (0.642,0.1212) -- (0.644,0.1247) -- (0.642,0.1283) -- (0.6379,0.1283) --
  (0.6358,0.1247) -- (0.6337,0.1283) -- (0.6296,0.1283) -- (0.6276,0.1247) -- (0.6296,0.1212) -- (0.6255,0.1212) --
  (0.6235,0.1176) -- (0.6255,0.114) -- (0.6296,0.114) -- (0.6276,0.1105) -- (0.6296,0.1069) -- (0.6255,0.1069) --
  (0.6235,0.1034) -- (0.6214,0.1069) -- (0.6173,0.1069) -- (0.6152,0.1034) -- (0.6173,0.0998) -- (0.6132,0.0998) --
  (0.6111,0.0962) -- (0.6132,0.0927) -- (0.6173,0.0927) -- (0.6152,0.0891) -- (0.6173,0.0855) -- (0.6214,0.0855) --
  (0.6235,0.0891) -- (0.6255,0.0855) -- (0.6296,0.0855) -- (0.6276,0.082) -- (0.6296,0.0784) -- (0.6255,0.0784) --
  (0.6235,0.0748) -- (0.6255,0.0713) -- (0.6296,0.0713) -- (0.6276,0.0677) -- (0.6296,0.0642) -- (0.6337,0.0642) --
  (0.6358,0.0677) -- (0.6379,0.0642) -- (0.642,0.0642) -- (0.644,0.0677) -- (0.642,0.0713) -- (0.6461,0.0713) --
  (0.6481,0.0748) -- (0.6502,0.0713) -- (0.6543,0.0713) -- (0.6523,0.0677) -- (0.6543,0.0642) -- (0.6584,0.0642) --
  (0.6605,0.0677) -- (0.6626,0.0642) -- (0.6667,0.0642) -- (0.6646,0.0606) -- (0.6667,0.057) -- (0.6626,0.057) --
  (0.6605,0.0535) -- (0.6626,0.0499) -- (0.6667,0.0499) -- (0.6646,0.0463) -- (0.6667,0.0428) -- (0.6626,0.0428) --
  (0.6605,0.0392) -- (0.6584,0.0428) -- (0.6543,0.0428) -- (0.6523,0.0392) -- (0.6543,0.0356) -- (0.6502,0.0356) --
  (0.6481,0.0321) -- (0.6502,0.0285) -- (0.6543,0.0285) -- (0.6523,0.0249) -- (0.6543,0.0214) -- (0.6584,0.0214) --
  (0.6605,0.0249) -- (0.6626,0.0214) -- (0.6667,0.0214) -- (0.6646,0.0178) -- (0.6667,0.0143) -- (0.6626,0.0143) --
  (0.6605,0.0107) -- (0.6626,0.0071) -- (0.6667,0.0071) -- (0.6646,0.0036) -- (0.6667,0);
\draw[trait4]
  (0.6667,0) -- (0.6708,0) -- (0.6728,0.0036) -- (0.6749,0) -- (0.679,0) -- (0.6811,0.0036) --
  (0.679,0.0071) -- (0.6831,0.0071) -- (0.6852,0.0107) -- (0.6872,0.0071) -- (0.6914,0.0071) -- (0.6893,0.0036) --
  (0.6914,0) -- (0.6955,0) -- (0.6975,0.0036) -- (0.6996,0) -- (0.7037,0) -- (0.7058,0.0036) --
  (0.7037,0.0071) -- (0.7078,0.0071) -- (0.7099,0.0107) -- (0.7078,0.0143) -- (0.7037,0.0143) -- (0.7058,0.0178) --
  (0.7037,0.0214) -- (0.7078,0.0214) -- (0.7099,0.0249) -- (0.7119,0.0214) -- (0.716,0.0214) -- (0.7181,0.0249) --
  (0.716,0.0285) -- (0.7202,0.0285) -- (0.7222,0.0321) -- (0.7243,0.0285) -- (0.7284,0.0285) -- (0.7263,0.0249) --
  (0.7284,0.0214) -- (0.7325,0.0214) -- (0.7346,0.0249) -- (0.7366,0.0214) -- (0.7407,0.0214) -- (0.7387,0.0178) --
  (0.7407,0.0143) -- (0.7366,0.0143) -- (0.7346,0.0107) -- (0.7366,0.0071) -- (0.7407,0.0071) -- (0.7387,0.0036) --
  (0.7407,0) -- (0.7449,0) -- (0.7469,0.0036) -- (0.749,0) -- (0.7531,0) -- (0.7551,0.0036) --
  (0.7531,0.0071) -- (0.7572,0.0071) -- (0.7593,0.0107) -- (0.7613,0.0071) -- (0.7654,0.0071) -- (0.7634,0.0036) --
  (0.7654,0) -- (0.7695,0) -- (0.7716,0.0036) -- (0.7737,0) -- (0.7778,0) -- (0.7798,0.0036) --
  (0.7778,0.0071) -- (0.7819,0.0071) -- (0.784,0.0107) -- (0.7819,0.0143) -- (0.7778,0.0143) -- (0.7798,0.0178) --
  (0.7778,0.0214) -- (0.7819,0.0214) -- (0.784,0.0249) -- (0.786,0.0214) -- (0.7901,0.0214) -- (0.7922,0.0249) --
  (0.7901,0.0285) -- (0.7942,0.0285) -- (0.7963,0.0321) -- (0.7942,0.0356) -- (0.7901,0.0356) -- (0.7922,0.0392) --
  (0.7901,0.0428) -- (0.786,0.0428) -- (0.784,0.0392) -- (0.7819,0.0428) -- (0.7778,0.0428) -- (0.7798,0.0463) --
  (0.7778,0.0499) -- (0.7819,0.0499) -- (0.784,0.0535) -- (0.7819,0.057) -- (0.7778,0.057) -- (0.7798,0.0606) --
  (0.7778,0.0642) -- (0.7819,0.0642) -- (0.784,0.0677) -- (0.786,0.0642) -- (0.7901,0.0642) -- (0.7922,0.0677) --
  (0.7901,0.0713) -- (0.7942,0.0713) -- (0.7963,0.0748) -- (0.7984,0.0713) -- (0.8025,0.0713) -- (0.8004,0.0677) --
  (0.8025,0.0642) -- (0.8066,0.0642) -- (0.8086,0.0677) -- (0.8107,0.0642) -- (0.8148,0.0642) -- (0.8169,0.0677) --
  (0.8148,0.0713) -- (0.8189,0.0713) -- (0.821,0.0748) -- (0.8189,0.0784) -- (0.8148,0.0784) -- (0.8169,0.082) --
  (0.8148,0.0855) -- (0.8189,0.0855) -- (0.821,0.0891) -- (0.823,0.0855) -- (0.8272,0.0855) -- (0.8292,0.0891) --
  (0.8272,0.0927) -- (0.8313,0.0927) -- (0.8333,0.0962) -- (0.8354,0.0927) -- (0.8395,0.0927) -- (0.8374,0.0891) --
  (0.8395,0.0855) -- (0.8436,0.0855) -- (0.8457,0.0891) -- (0.8477,0.0855) -- (0.8519,0.0855) -- (0.8498,0.082) --
  (0.8519,0.0784) -- (0.8477,0.0784) -- (0.8457,0.0748) -- (0.8477,0.0713) -- (0.8519,0.0713) -- (0.8498,0.0677) --
  (0.8519,0.0642) -- (0.856,0.0642) -- (0.858,0.0677) -- (0.8601,0.0642) -- (0.8642,0.0642) -- (0.8663,0.0677) --
  (0.8642,0.0713) -- (0.8683,0.0713) -- (0.8704,0.0748) -- (0.8724,0.0713) -- (0.8765,0.0713) -- (0.8745,0.0677) --
  (0.8765,0.0642) -- (0.8807,0.0642) -- (0.8827,0.0677) -- (0.8848,0.0642) -- (0.8889,0.0642) -- (0.8868,0.0606) --
  (0.8889,0.057) -- (0.8848,0.057) -- (0.8827,0.0535) -- (0.8848,0.0499) -- (0.8889,0.0499) -- (0.8868,0.0463) --
  (0.8889,0.0428) -- (0.8848,0.0428) -- (0.8827,0.0392) -- (0.8807,0.0428) -- (0.8765,0.0428) -- (0.8745,0.0392) --
  (0.8765,0.0356) -- (0.8724,0.0356) -- (0.8704,0.0321) -- (0.8724,0.0285) -- (0.8765,0.0285) -- (0.8745,0.0249) --
  (0.8765,0.0214) -- (0.8807,0.0214) -- (0.8827,0.0249) -- (0.8848,0.0214) -- (0.8889,0.0214) -- (0.8868,0.0178) --
  (0.8889,0.0143) -- (0.8848,0.0143) -- (0.8827,0.0107) -- (0.8848,0.0071) -- (0.8889,0.0071) -- (0.8868,0.0036) --
  (0.8889,0) -- (0.893,0) -- (0.8951,0.0036) -- (0.8971,0) -- (0.9012,0) -- (0.9033,0.0036) --
  (0.9012,0.0071) -- (0.9053,0.0071) -- (0.9074,0.0107) -- (0.9095,0.0071) -- (0.9136,0.0071) -- (0.9115,0.0036) --
  (0.9136,0) -- (0.9177,0) -- (0.9198,0.0036) -- (0.9218,0) -- (0.9259,0) -- (0.928,0.0036) --
  (0.9259,0.0071) -- (0.93,0.0071) -- (0.9321,0.0107) -- (0.93,0.0143) -- (0.9259,0.0143) -- (0.928,0.0178) --
  (0.9259,0.0214) -- (0.93,0.0214) -- (0.9321,0.0249) -- (0.9342,0.0214) -- (0.9383,0.0214) -- (0.9403,0.0249) --
  (0.9383,0.0285) -- (0.9424,0.0285) -- (0.9444,0.0321) -- (0.9465,0.0285) -- (0.9506,0.0285) -- (0.9486,0.0249) --
  (0.9506,0.0214) -- (0.9547,0.0214) -- (0.9568,0.0249) -- (0.9588,0.0214) -- (0.963,0.0214) -- (0.9609,0.0178) --
  (0.963,0.0143) -- (0.9588,0.0143) -- (0.9568,0.0107) -- (0.9588,0.0071) -- (0.963,0.0071) -- (0.9609,0.0036) --
  (0.963,0) -- (0.9671,0) -- (0.9691,0.0036) -- (0.9712,0) -- (0.9753,0) -- (0.9774,0.0036) --
  (0.9753,0.0071) -- (0.9794,0.0071) -- (0.9815,0.0107) -- (0.9835,0.0071) -- (0.9877,0.0071) -- (0.9856,0.0036) --
  (0.9877,0) -- (0.9918,0) -- (0.9938,0.0036) -- (0.9959,0) -- (1,0);

        \node[etiq, firebrick3, below] at (0.1667,-0.03) {$w_1(K)$};
        \node[etiq, steelblue4, above left] at (0.36,0.25) {$w_2(K)$};
        \node[etiq, coveraccent!80!black, above right] at (0.64,0.25) {$w_3(K)$};
        \node[etiq, vertpiece, below] at (0.8333,-0.03) {$w_4(K)$};
      \end{scope}
    \end{tikzpicture}}
  \small Tapis : huit homothéties de rapport $\tfrac13$. Cantor : deux
  copies, $\tfrac13x$ et $\tfrac13x+\tfrac23$.
\end{frame}

\begin{frame}{Remarque 3.2 et pour la classe 3.2}
  \begin{Rem}{Remarque 3.2 --- Une autosimilarité parfois seulement statistique}
    \footnotesize Une côte rocheuse, un chou-fleur, un réseau de rivières
    ne sont jamais exactement autosimilaires : un morceau agrandi ne
    reproduit pas le tout, mais en a les caractéristiques statistiques,
    et seulement sur une plage d'échelles (pas au grain de sable). C'est
    l'autosimilarité statistique, au cœur des travaux de Mandelbrot.
  \end{Rem}
  \begin{Classe}{Une figure semblable à une partie d'elle-même}
    \footnotesize Une figure autosimilaire est une figure semblable à
    chacune de ses parties. Les élèves joignent les milieux des côtés
    d'un triangle, mesurent, constatent quatre petits triangles
    isométriques semblables au grand, de rapport $\tfrac12$. Recommencer
    dans chaque petit triangle en noircissant celui du centre amorce la
    construction de Sierpiński.
  \end{Classe}
\end{frame}
